\documentclass[fs12,seriesSau]{korfarm-base}
\usepackage{korfarm-saussure}

% ============================================================
% passage-note.tex (v11)
% 지문 박스 + 첫 페이지 우측 메모란 (동적 높이)
% 정책: PAGE_BREAK_POLICY.md §3, §4 참조
%
% 변경 (v11):
%  · 지문 박스 폭 확대: enlarge right by=-3.7cm  (메모 3.4cm + gap 3mm)
%  · 메모란 동적 높이: frame.north east ~ frame.south east 사이 자동 채움
%  · 메모 줄선: clip 영역으로 박스 내부만 채움 (박스 높이 모르고도 작동)
%  · 문단 들여쓰기 활성화: tcolorbox 내부 \parindent=1em, \noindent 제거
%
% 사용:
%   \kfPassageNote{제목}{지문 본문}      → 메모란 포함
%   \kfPassageNoteNT{지문 본문}          → 메모란 포함, 제목 없음
%   \kfPassageFull{제목}{본문}           → 메모란 없음 (활동 내 보기)
% ============================================================

\makeatletter

% 메모란 규격 (cm 고정)
\newlength{\kf@memoW}
\setlength{\kf@memoW}{3.4cm}
\newlength{\kf@memoGap}
\setlength{\kf@memoGap}{3mm}

% --- 메인: 제목 있는 버전 ---
\newcommand{\kf@passagenote@with}[2]{%
  \par\ifkfPassageNewPage\ifdim\pagetotal>0.4\textheight\clearpage\fi\fi\smallskip\needspace{6\baselineskip}%
  \begin{tcolorbox}[
    enhanced, breakable,
    width=\dimexpr\linewidth-4cm\relax,
    colback=kfPassageBg, colframe=kfPrimary!70,
    boxrule=0.6pt, arc=3pt,
    left=\kfBoxPad, right=\kfBoxPad, top=\dimexpr\kfBoxPad+8pt\relax, bottom=\kfBoxPad,
    title={\small\bfseries\color{kfPrimary} #1},
    coltitle=kfPrimary,
    colbacktitle=kfPrimaryLight!60,
    attach boxed title to top left={yshift=-2.4mm, xshift=8pt},
    boxed title style={colframe=kfPrimary!50, boxrule=0.5pt, arc=2pt},
    parbox=false,
    overlay unbroken and first={%
      \begin{scope}
        \draw[kfRule, line width=0.4pt, fill=kfNoteBg, rounded corners=2pt]
          ([xshift=\kf@memoGap]frame.north east) rectangle
          ([xshift=\kf@memoGap+\kf@memoW]frame.south east);
        \node[anchor=north east, font=\footnotesize\bfseries\color{kfMute}, inner sep=3pt]
          at ([xshift=\kf@memoGap+\kf@memoW, yshift=-2pt]frame.north east) {메모};
        \node[anchor=north east, fill=kfPrimary!15, text=kfPrimary,
              font=\scriptsize\bfseries, rounded corners=2pt,
              inner xsep=4pt, inner ysep=2pt]
          at ([xshift=\kf@memoGap+\kf@memoW-3pt, yshift=-19pt]frame.north east) {\kf@memocat};
        \begin{scope}
          \path[clip]
            ([xshift=\kf@memoGap]frame.north east) rectangle
            ([xshift=\kf@memoGap+\kf@memoW]frame.south east);
          \foreach \i in {1,...,40}{%
            \draw[kfRule, line width=0.25pt]
              ([xshift=\kf@memoGap+2mm, yshift=-7mm - \i*5mm]frame.north east) --
              ([xshift=\kf@memoGap+\kf@memoW-2mm, yshift=-7mm - \i*5mm]frame.north east);%
          }
        \end{scope}
      \end{scope}
    }
  ]
    \setlength{\parindent}{1em}%
    \hspace*{1em}\ignorespaces#2%
  \end{tcolorbox}\par\smallskip
}

% --- 제목 없는 버전 ---
\newcommand{\kf@passagenote@notitle}[1]{%
  \par\ifkfPassageNewPage\ifdim\pagetotal>0.4\textheight\clearpage\fi\fi\smallskip\needspace{6\baselineskip}%
  \begin{tcolorbox}[
    enhanced, breakable,
    width=\dimexpr\linewidth-4cm\relax,
    colback=kfPassageBg, colframe=kfPrimary!70,
    boxrule=0.6pt, arc=3pt,
    left=\kfBoxPad, right=\kfBoxPad, top=\dimexpr\kfBoxPad+8pt\relax, bottom=\kfBoxPad,
    parbox=false,
    overlay unbroken and first={%
      \begin{scope}
        \draw[kfRule, line width=0.4pt, fill=kfNoteBg, rounded corners=2pt]
          ([xshift=\kf@memoGap]frame.north east) rectangle
          ([xshift=\kf@memoGap+\kf@memoW]frame.south east);
        \node[anchor=north east, font=\footnotesize\bfseries\color{kfMute}, inner sep=3pt]
          at ([xshift=\kf@memoGap+\kf@memoW, yshift=-2pt]frame.north east) {메모};
        \node[anchor=north east, fill=kfPrimary!15, text=kfPrimary,
              font=\scriptsize\bfseries, rounded corners=2pt,
              inner xsep=4pt, inner ysep=2pt]
          at ([xshift=\kf@memoGap+\kf@memoW-3pt, yshift=-19pt]frame.north east) {\kf@memocat};
        \begin{scope}
          \path[clip]
            ([xshift=\kf@memoGap]frame.north east) rectangle
            ([xshift=\kf@memoGap+\kf@memoW]frame.south east);
          \foreach \i in {1,...,40}{%
            \draw[kfRule, line width=0.25pt]
              ([xshift=\kf@memoGap+2mm, yshift=-7mm - \i*5mm]frame.north east) --
              ([xshift=\kf@memoGap+\kf@memoW-2mm, yshift=-7mm - \i*5mm]frame.north east);%
          }
        \end{scope}
      \end{scope}
    }
  ]
    \setlength{\parindent}{1em}%
    \hspace*{1em}\ignorespaces#1%
  \end{tcolorbox}\par\smallskip
}

\newcommand{\kfPassageNote}[2]{%
  \def\kf@tmptitle{#1}%
  \ifx\kf@tmptitle\@empty
    \kf@passagenote@notitle{#2}%
  \else
    \kf@passagenote@with{#1}{#2}%
  \fi
}
\newcommand{\kfPassageNoteNT}[1]{\kf@passagenote@notitle{#1}}
\makeatother

% ============================================================
% 풀폭 지문 박스 (메모란 없음) — 활동 내 보기 박스
% PAGE_BREAK_POLICY.md §3
% ============================================================
\makeatletter
\newcommand{\kf@passagefull@with}[2]{%
  \par\ifkfPassageNewPage\ifdim\pagetotal>0.4\textheight\clearpage\fi\fi\smallskip\needspace{6\baselineskip}%
  \begin{tcolorbox}[
    enhanced, breakable,
    colback=kfPassageBg, colframe=kfPrimary!70,
    boxrule=0.6pt, arc=3pt,
    left=\kfBoxPad, right=\kfBoxPad, top=\dimexpr\kfBoxPad+8pt\relax, bottom=\kfBoxPad,
    title={\small\bfseries\color{kfPrimary} #1},
    coltitle=kfPrimary,
    colbacktitle=kfPrimaryLight!60,
    attach boxed title to top left={yshift=-2.4mm, xshift=8pt},
    boxed title style={colframe=kfPrimary!50, boxrule=0.5pt, arc=2pt},
    parbox=false
  ]
    \setlength{\parindent}{1em}%
    \hspace*{1em}\ignorespaces#2%
  \end{tcolorbox}\par\smallskip
}
\newcommand{\kf@passagefull@notitle}[1]{%
  \par\ifkfPassageNewPage\ifdim\pagetotal>0.4\textheight\clearpage\fi\fi\smallskip\needspace{6\baselineskip}%
  \begin{tcolorbox}[
    enhanced, breakable,
    colback=kfPassageBg, colframe=kfPrimary!70,
    boxrule=0.6pt, arc=3pt,
    left=\kfBoxPad, right=\kfBoxPad, top=\dimexpr\kfBoxPad+8pt\relax, bottom=\kfBoxPad,
    parbox=false
  ]
    \setlength{\parindent}{1em}%
    \hspace*{1em}\ignorespaces#1%
  \end{tcolorbox}\par\smallskip
}
\newcommand{\kfPassageFull}[2]{%
  \def\kf@tmptitle{#1}%
  \ifx\kf@tmptitle\@empty
    \kf@passagefull@notitle{#2}%
  \else
    \kf@passagefull@with{#1}{#2}%
  \fi
}
\newcommand{\kfPassageFullNT}[1]{\kf@passagefull@notitle{#1}}
\makeatother

% ============================================================
% 작은 문장 박스 (문장 독해용) — PAGE_BREAK_POLICY.md §2.5
% ============================================================
\newcommand{\kfSentenceBox}[2]{%
  \par\smallskip\needspace{3\baselineskip}%
  \noindent\begin{tcolorbox}[
    enhanced, breakable=false,
    colback=kfPassageBg, colframe=kfMute,
    boxrule=0.4pt, arc=2pt,
    left=8pt, right=6pt, top=4pt, bottom=4pt,
    title={\footnotesize\bfseries\color{kfPrimary} #1},
    coltitle=kfPrimary, colbacktitle=kfPaper,
    attach boxed title to top left={yshift=-1.6mm, xshift=4pt},
    boxed title style={colframe=kfMute, boxrule=0.3pt, arc=1pt}
  ]%
    \setstretch{1.3}#2%
  \end{tcolorbox}\par\smallskip%
}

% ============================================================
% 지문 본문 아래 안내/정리 박스 (v16 신규)
% 사용: \kfPassageHint{한 줄 정리}{본문 안내 텍스트}
% ============================================================
\newcommand{\kfPassageHint}[2]{%
  \par\smallskip\needspace{4\baselineskip}%
  \begin{tcolorbox}[
    enhanced, breakable=false,
    colback=kfPrimary!5, colframe=kfPrimary!40,
    boxrule=0.4pt, arc=2pt,
    left=8pt, right=6pt, top=4pt, bottom=4pt,
    title={\footnotesize\bfseries\color{kfPrimary} #1},
    coltitle=kfPrimary, colbacktitle=kfPaper,
    attach boxed title to top left={yshift=-1.5mm, xshift=4pt},
    boxed title style={colframe=kfPrimary!40, boxrule=0.3pt, arc=1pt}
  ]%
    \small\setstretch{1.3}#2%
  \end{tcolorbox}\par\smallskip%
}

% ============================================================
% 환경 버전 (호환성) — 메모란 포함
% ============================================================
\makeatletter
\newenvironment{kfPassageNoteEnv}[1]%
  {\par\needspace{6\baselineskip}%
   \def\kf@psrc{#1}%
   \begin{tcolorbox}[
     enhanced, breakable,
     width=\dimexpr\linewidth-4cm\relax,
     colback=kfPassageBg, colframe=kfPrimary!70,
     boxrule=0.6pt, arc=3pt,
     left=\kfBoxPad, right=\kfBoxPad, top=\dimexpr\kfBoxPad+8pt\relax, bottom=\kfBoxPad,
     title={\small\bfseries\color{kfPrimary} \kf@psrc},
     coltitle=kfPrimary, colbacktitle=kfPrimaryLight!60,
     attach boxed title to top left={yshift=-2.4mm, xshift=8pt},
     boxed title style={colframe=kfPrimary!50, boxrule=0.5pt, arc=2pt},
     parbox=false,
     overlay unbroken and first={%
       \begin{scope}
         \draw[kfRule, line width=0.4pt, fill=kfNoteBg, rounded corners=2pt]
           ([xshift=\kf@memoGap]frame.north east) rectangle
           ([xshift=\kf@memoGap+\kf@memoW]frame.south east);
         \node[anchor=north east, font=\footnotesize\bfseries\color{kfMute}, inner sep=3pt]
           at ([xshift=\kf@memoGap+\kf@memoW, yshift=-2pt]frame.north east) {메모};
         \begin{scope}
           \path[clip]
             ([xshift=\kf@memoGap]frame.north east) rectangle
             ([xshift=\kf@memoGap+\kf@memoW]frame.south east);
           \foreach \i in {1,...,40}{%
             \draw[kfRule, line width=0.25pt]
               ([xshift=\kf@memoGap+2mm, yshift=-7mm - \i*5mm]frame.north east) --
               ([xshift=\kf@memoGap+\kf@memoW-2mm, yshift=-7mm - \i*5mm]frame.north east);%
           }
         \end{scope}
       \end{scope}
     }
   ]%
   \setlength{\parindent}{1em}%
  }%
  {\end{tcolorbox}\par\smallskip}
\makeatother

% ============================================================
% condition-box.tex
% <조건> 박스 컴포넌트 (tcolorbox)
% 사용:
%   \begin{kfCondition}
%     \item 첫째 조건
%     \item 둘째 조건
%   \end{kfCondition}
% ============================================================

\newenvironment{kfCondition}%
  {\par\smallskip\needspace{4\baselineskip}%
   \begin{tcolorbox}[
     enhanced, breakable=false,
     colback=kfPrimaryLight!40, colframe=kfPrimary,
     boxrule=0.5pt, arc=2pt,
     left=\kfBoxPad, right=\kfBoxPad, top=\kfBoxPad, bottom=\kfBoxPad,
     title={\small\bfseries\color{kfPrimary} \,조건\,},
     coltitle=kfPrimary, colbacktitle=kfPaper,
     attach boxed title to top left={yshift=-2mm, xshift=6pt},
     boxed title style={colframe=kfPrimary, boxrule=0.4pt, arc=1pt}
   ]%
   \begin{enumerate}[leftmargin=16pt, itemsep=1pt, topsep=1pt,
                    label=\textbf{\arabic*.}]}%
  {\end{enumerate}\end{tcolorbox}\par\smallskip}

% ============================================================
% evidence-box.tex
% <보기> 박스 컴포넌트
% 사용:
%   \begin{kfEvidence}
%     보기 본문 ...
%   \end{kfEvidence}
% ============================================================

\newenvironment{kfEvidence}%
  {\par\smallskip\needspace{4\baselineskip}%
   \begin{tcolorbox}[
     enhanced, breakable=false,
     colback=kfPaper, colframe=kfMute,
     boxrule=0.5pt, arc=2pt,
     left=\kfBoxPad, right=\kfBoxPad, top=\kfBoxPad, bottom=\kfBoxPad,
     title={\small\bfseries\color{kfInk} \,보기\,},
     coltitle=kfInk, colbacktitle=kfPrimaryLight!50,
     attach boxed title to top left={yshift=-2mm, xshift=6pt},
     boxed title style={colframe=kfMute, boxrule=0.4pt, arc=1pt}
   ]%
   \leavevmode\mbox{}\ignorespaces%
  }%
  {\end{tcolorbox}\par\smallskip}

% ============================================================
% choice-list.tex
% 5선지 (① ② ③ ④ ⑤) 자동 들여쓰기 컴포넌트
% 사용:
%   \begin{kfChoices}
%     \kfChoice 첫째 선택지
%     \kfChoice 둘째 선택지
%     \kfChoice 셋째 선택지
%     \kfChoice 넷째 선택지
%     \kfChoice 다섯째 선택지
%   \end{kfChoices}
% ============================================================

\newcounter{kfChoiceCnt}
\newcommand{\kfChoiceMark}[1]{%
  \ifcase#1\relax%
  \or ①\or ②\or ③\or ④\or ⑤\or ⑥\or ⑦\or ⑧\or ⑨%
  \fi
}

% 환경: 자동 번호
\newenvironment{kfChoices}%
  {\setcounter{kfChoiceCnt}{0}%
   \par\smallskip%
   \begin{list}{}{%
     \setlength{\leftmargin}{20pt}%
     \setlength{\itemindent}{0pt}%
     \setlength{\labelwidth}{18pt}%
     \setlength{\labelsep}{6pt}%
     \setlength{\itemsep}{2pt}%
     \setlength{\topsep}{1pt}%
     \setlength{\parsep}{0pt}%
     \renewcommand{\makelabel}[1]{##1\hss}%
   }}%
  {\end{list}\par\smallskip}

\newcommand{\kfChoice}{%
  \stepcounter{kfChoiceCnt}%
  \item[\textbf{\kfChoiceMark{\value{kfChoiceCnt}}}]\ignorespaces%
}

% --- 한 줄 5선지 (짧은 선택지용) ---
\newcommand{\kfChoicesInline}[5]{%
  \par\smallskip%
  \noindent\textbf{①}~#1\quad\textbf{②}~#2\quad\textbf{③}~#3\quad\textbf{④}~#4\quad\textbf{⑤}~#5%
  \par\smallskip%
}

% ============================================================
% answer-note.tex
% 답안 작성란 (노트 스타일, 줄친 영역)
% 사용:
%   \kfAnswerLines{5}        % 5줄 답안란
%   \kfAnswerNote{서술형 답안란}{8} % 라벨+8줄
% ============================================================

% 단일 줄 (필요 시 인라인)
\newcommand{\kfAnswerLine}{%
  \par\noindent\textcolor{kfRule}{\rule{\linewidth}{0.4pt}}\par
}

% N줄 답안란 (간단)
\newcommand{\kfAnswerLines}[1]{%
  \par\smallskip%
  \begin{minipage}{\linewidth}%
  \foreach \i in {1,...,#1}{%
    \textcolor{kfRule}{\rule{\linewidth}{0.4pt}}\\[6pt]%
  }%
  \end{minipage}%
  \par\smallskip%
}

% 라벨이 있는 노트형 답안란
\newcommand{\kfAnswerNote}[2]{%
  \par\smallskip\needspace{#2\baselineskip}%
  \begin{tcolorbox}[
    enhanced, breakable=false,
    colback=kfPaper, colframe=kfRule,
    boxrule=0.4pt, arc=2pt,
    left=\kfBoxPad, right=\kfBoxPad, top=\kfBoxPad, bottom=\kfBoxPad,
    title={\footnotesize\bfseries\color{kfMute} #1},
    coltitle=kfMute, colbacktitle=kfPaper,
    attach boxed title to top left={yshift=-1.5mm, xshift=6pt},
    boxed title style={colframe=kfRule, boxrule=0.3pt, arc=1pt}
  ]%
  \foreach \i in {1,...,#2}{%
    \textcolor{kfRule}{\rule{\linewidth}{0.3pt}}\\[6pt]%
  }%
  \end{tcolorbox}\par\smallskip%
}

% 짧은 빈칸 (단답형)
\newcommand{\kfBlank}[1][3cm]{\underline{\hspace{#1}}}
% 넓은 빈칸 (서술 한 줄)
\newcommand{\kfBlankWide}[1][6cm]{\underline{\hspace{#1}}}
% 괄호 빈칸
\newcommand{\kfBlankParen}{(\,\hspace{1cm}\,)}
% O/X 빈칸
\newcommand{\kfBlankOX}{(\,\hspace{1.5cm}\,)}

% ============================================================
% question-block.tex
% 문제 블록 (번호 배지 + 발문 + 보기/조건 + 선택지 + 답안란)
% 모든 구성을 하나의 minipage로 묶어 단/페이지 단절 금지
%
% 사용:
%   \begin{kfQBlock}{1}{객관식}
%     \kfQStem{다음 글에 대한 설명으로 가장 적절한 것은?}
%     \begin{kfEvidence} ... \end{kfEvidence}    % 선택
%     \begin{kfCondition} ... \end{kfCondition}  % 선택
%     \begin{kfChoices}
%       \kfChoice ...
%     \end{kfChoices}
%     \kfAnswerLines{2}                            % 선택
%   \end{kfQBlock}
% ============================================================

% 무단절 minipage로 묶고, 발문/보기/조건/선택지/답안란을 자유롭게 배치
% 사용자 피드백 #4: [객관식]/[서술형] 등 텍스트 라벨 제거 — 번호 배지만 표시
% #2(유형 라벨)는 호환성 인자로만 받고 출력하지 않음.
\newenvironment{kfQBlock}[2]%
  {\par\smallskip\needspace{6\baselineskip}%
   \noindent\begin{minipage}{\linewidth}%
   \samepage%
   \par\noindent
   \kfQNum{#1}\hspace{4pt}%
   \begin{minipage}[t]{\dimexpr\linewidth-30pt\relax}%
  }%
  {\end{minipage}\end{minipage}\par\smallskip}

% 문제 발문
\newcommand{\kfQStem}[1]{%
  \par\noindent#1\par\vspace{2pt}%
}

% 짧은 소문항 라벨 (1), (2), (3)
\newcounter{kfSubQ}
\newcommand{\kfSubQReset}{\setcounter{kfSubQ}{0}}
\newcommand{\kfSubQ}{%
  \stepcounter{kfSubQ}%
  \par\noindent\textbf{(\arabic{kfSubQ})}~\ignorespaces%
}

% ============================================================
% activity-table.tex
% 활동: 정리표 / 내용 확인표 (마크다운 표 → tabularx 변환)
% 사용:
%   % 일반 tabular (직접 컬럼 명시)
%   \begin{kfActTable}{|l|X|X|}
%     \kfTHead{항목} & \kfTHead{설명} & \kfTHead{학생 작성} \\\hline
%     ① & ... & \kfTBlank \\\hline
%   \end{kfActTable}
%
%   % adjustbox 자동 축소 (정해진 폭으로 강제)
%   \kfActTableWrap{
%     ... (\begin{tabular}...\end{tabular} 코드) ...
%   }
% ============================================================

% 사용 시 본문 마지막 행은 \\\hline 으로 끝나야 함.
% v17.1: 흰색 mask로 Canva 배경 본문 침범 차단
\newenvironment{kfActTable}[1]%
  {\par\smallskip\needspace{4\baselineskip}%
   \renewcommand{\arraystretch}{1.3}%
   \arrayrulecolor{kfRule}%
   \begin{tcolorbox}[blanker, colback=white, left=2pt, right=2pt, top=2pt, bottom=2pt]%
   \noindent\begin{adjustbox}{max width=\linewidth}%
   \begin{tabular}{#1}%
   \hline}%
  {\end{tabular}%
   \end{adjustbox}%
   \end{tcolorbox}%
   \arrayrulecolor{black}%
   \par\smallskip}

% 헤더 행 헬퍼
\newcommand{\kfTHead}[1]{\textbf{\color{kfPrimary} #1}}
\newcommand{\kfTHeadRow}[1]{\rowcolor{kfPrimaryLight!50}\textbf{\color{kfPrimary} #1}}

% 빈 셀 (학생 작성용)
\newcommand{\kfTBlank}{\rule[-1.6em]{0pt}{3.4em}}
\newcommand{\kfTBlankSm}{\rule[-1.0em]{0pt}{2.4em}}

% adjustbox 강제 축소 래퍼 — Python에서 변환된 raw tabular 블록을 감쌀 때 사용
\newcommand{\kfActTableWrap}[1]{%
  \par\smallskip\needspace{4\baselineskip}%
  \begin{tcolorbox}[blanker, colback=white, left=2pt, right=2pt, top=2pt, bottom=2pt]%
  \noindent\begin{adjustbox}{max width=\linewidth, center}%
  #1%
  \end{adjustbox}%
  \end{tcolorbox}\par\smallskip%
}

% ============================================================
% activity-vocab.tex
% 어휘 학습 표 (3열: 번호 - 뜻 - 빈칸)
% 사용:
%   \begin{kfVocab}
%     \kfVocabItem{1}{눈으로 보고 마음으로 깨달음}
%     \kfVocabItem{2}{사물의 이치를 따져 헤아림}
%   \end{kfVocab}
% ============================================================

% adjustbox로 폭 강제 + 일반 tabular + p{} 열 사용
\newenvironment{kfVocab}%
  {\par\smallskip\needspace{4\baselineskip}%
   \renewcommand{\arraystretch}{1.35}%
   \arrayrulecolor{kfRule}%
   \noindent\begin{adjustbox}{max width=\linewidth, center}%
   \begin{tabular}{|>{\centering\arraybackslash}m{1.0cm}|m{8.5cm}|>{\centering\arraybackslash}m{4.0cm}|}%
   \hline
   \rowcolor{kfPrimaryLight!50}%
   \multicolumn{1}{|c|}{\textbf{\color{kfPrimary}번호}} &
   \multicolumn{1}{c|}{\textbf{\color{kfPrimary}뜻풀이}} &
   \multicolumn{1}{c|}{\textbf{\color{kfPrimary}어휘 (정답 작성)}} \\\hline
   \ignorespaces}%
  {\end{tabular}%
   \end{adjustbox}%
   \arrayrulecolor{black}%
   \par\smallskip}

\newcommand{\kfVocabItem}[2]{%
  \textbf{#1} & #2 & \rule[-1.0em]{0pt}{2.4em}\\\hline%
}

% ============================================================
% activity-blank.tex
% 빈칸 채우기 활동
% 사용:
%   \begin{kfBlankList}
%     \kfBlankItem 첫 번째 문장에서 (\kfBlank)은(는) ...
%     \kfBlankItem 둘째 문장에서 ...
%   \end{kfBlankList}
% ============================================================

\newenvironment{kfBlankList}%
  {\par\smallskip%
   \begin{enumerate}[leftmargin=20pt, itemsep=4pt, topsep=2pt,
                    label=\textbf{\arabic*.}, labelsep=6pt]}%
  {\end{enumerate}\par\smallskip}

\newcommand{\kfBlankItem}{\item}

% 텍스트 안에 박스형 빈칸 (단어 1개 채우기)
\newcommand{\kfBlankBox}[1][2.4cm]{%
  \tikz[baseline=-0.5ex]{%
    \draw[kfRule, line width=0.4pt] (0,0) -- ++(#1,0);%
  }%
}

% ============================================================
% activity-ox.tex (v15)
% O/X 표기 활동 — 그래픽 디자인
% 사용:
%   \begin{kfOXList}
%     \kfOXItem 글쓴이는 자연을 예찬하고 있다.\kfOX
%     \kfOXItem 1연과 4연은 수미상관 구조이다.\kfOX
%   \end{kfOXList}
%
% \kfOX = 우측에 [O] [X] 두 박스 (학생이 하나 동그라미)
% ============================================================

\newenvironment{kfOXList}%
  {\par\smallskip%
   \begin{enumerate}[leftmargin=20pt, itemsep=5pt, topsep=2pt,
                    label=\textbf{\arabic*.}, labelsep=6pt]}%
  {\end{enumerate}\par\smallskip}

\newcommand{\kfOXItem}{\item}

% v16: O/X 그래픽 박스 키움 (18×14 → 24×20pt) + 영역색 강조
\newcommand{\kfOX}{%
  \hfill\nolinebreak
  \tikz[baseline=-0.9ex]{%
    \node[draw=kfPrimary, fill=kfPrimary!5, line width=0.6pt, rounded corners=3pt,
          minimum width=24pt, minimum height=20pt, inner sep=0pt,
          font=\normalsize\bfseries\color{kfPrimary}] (ob) {O};%
    \node[draw=kfMute, fill=kfMute!5, line width=0.6pt, rounded corners=3pt,
          minimum width=24pt, minimum height=20pt, inner sep=0pt,
          font=\normalsize\bfseries\color{kfMute},
          right=4pt of ob] (xb) {X};%
  }%
}

% ============================================================
% activity-connect.tex
% 좌우 선 긋기 활동 (2단 양식 + 중앙 여백)
% 사용:
%   \begin{kfConnect}
%     \kfPair{사르트르}{실존은 본질에 앞선다}
%     \kfPair{니체}{초인 사상}
%     \kfPair{하이데거}{현존재}
%   \end{kfConnect}
% ============================================================

\newenvironment{kfConnect}%
  {\par\smallskip%
   \renewcommand{\arraystretch}{1.6}%
   \begin{tabular}{@{}>{\raggedleft\arraybackslash}m{4.5cm}
                     @{\hspace{2.5cm}}
                     >{\raggedright\arraybackslash}m{6.5cm}@{}}}%
  {\end{tabular}\par\smallskip}

\newcommand{\kfPair}[2]{%
  \tikz[baseline]{\node[draw=kfRule, fill=kfPrimaryLight!30, rounded corners=2pt,
        inner sep=4pt, font=\small]{#1};} \tikz[baseline]{\node[circle, draw=kfMute, fill=kfPaper, inner sep=1.5pt]{};}
  &
  \tikz[baseline]{\node[circle, draw=kfMute, fill=kfPaper, inner sep=1.5pt]{};} \tikz[baseline]{\node[draw=kfRule, fill=kfNoteBg, rounded corners=2pt,
        inner sep=4pt, font=\small]{#2};} \\
}

% ============================================================
% activity-writing.tex
% 글쓰기 활동 (큰 노트 영역)
% 사용:
%   \kfWriting{독후감}{12}     % 라벨 + 12줄 큰 노트
%   \kfWritingPrompt{프롬프트 텍스트}
% ============================================================

\newcommand{\kfWritingPrompt}[1]{%
  \par\smallskip%
  \begin{tcolorbox}[
    enhanced, breakable=false,
    colback=kfPrimaryLight!30, colframe=kfPrimary,
    boxrule=0.4pt, arc=3pt,
    left=\kfBoxPad, right=\kfBoxPad, top=\kfBoxPad, bottom=\kfBoxPad,
  ]%
  \small\textbf{\color{kfPrimary}글쓰기 안내}\\[2pt]%
  #1
  \end{tcolorbox}\par\smallskip%
}

\newcommand{\kfWriting}[2]{%
  \par\smallskip\needspace{\numexpr#2+2\relax\baselineskip}%
  \begin{tcolorbox}[
    enhanced, breakable=false,
    colback=kfPaper, colframe=kfRule,
    boxrule=0.4pt, arc=2pt,
    left=\kfBoxPad, right=\kfBoxPad, top=\kfBoxPad, bottom=\kfBoxPad,
    title={\footnotesize\bfseries\color{kfMute} #1},
    coltitle=kfMute, colbacktitle=kfPaper,
    attach boxed title to top left={yshift=-1.5mm, xshift=6pt},
    boxed title style={colframe=kfRule, boxrule=0.3pt, arc=1pt}
  ]%
  \foreach \i in {1,...,#2}{%
    \textcolor{kfRule}{\rule{\linewidth}{0.3pt}}\\[7pt]%
  }%
  \end{tcolorbox}\par\smallskip%
}

% ============================================================
% activity-structure.tex
% 구조도 (트리 + 빈칸)
% 사용:
%   \begin{kfStructure}
%     \kfNode{0}{서론}{문제 제기}
%     \kfNode{1}{본론}{(\,\quad\,)}
%     \kfNode{1}{본론}{근거 2}
%     \kfNode{0}{결론}{주장 강화}
%   \end{kfStructure}
% ============================================================

\newenvironment{kfStructure}%
  {\par\smallskip\needspace{6\baselineskip}%
   \begin{tcolorbox}[
     enhanced, breakable=false,
     colback=kfPaper, colframe=kfRule,
     boxrule=0.4pt, arc=2pt,
     left=\kfBoxPad, right=\kfBoxPad, top=\kfBoxPad, bottom=\kfBoxPad,
     title={\footnotesize\bfseries\color{kfMute} 글의 구조},
     coltitle=kfMute, colbacktitle=kfPrimaryLight!40,
     attach boxed title to top left={yshift=-1.5mm, xshift=6pt},
     boxed title style={colframe=kfRule, boxrule=0.3pt, arc=1pt}
   ]%
   \begin{tabbing}
     \hspace{1.4cm}\=\hspace{1.4cm}\=\hspace{8cm}\kill}%
  {\end{tabbing}\end{tcolorbox}\par\smallskip}

% 노드: #1=들여쓰기 단계 (0~3), #2=라벨, #3=내용
\newcommand{\kfNode}[3]{%
  \ifcase#1\relax%
    \>\>\textbf{\color{kfPrimary} ▶ #2}\quad #3\\
  \or
    \>\>\quad\textbf{\color{kfMute} ◦ #2}\quad #3\\
  \or
    \>\>\quad\quad\textbf{\color{kfMute} - #2}\quad #3\\
  \or
    \>\>\quad\quad\quad\textbf{\color{kfMute} \textperiodcentered\ #2}\quad #3\\
  \fi
}

% 빈칸 노드
\newcommand{\kfNodeBlank}[2]{\kfNode{#1}{#2}{(\hspace{2.5cm})}}

% ============================================================
% activity-sentence-reading.tex
% 문장 독해 활동 — 문장 박스 1개 + 그에 묶인 문제 N개를 한 minipage로
% 사용:
%   \begin{kfSentenceUnit}{1}{"바람은 내 귀에 속삭이며..."}
%     \kfSentQ{(1)}{서술어 '속삭이며'의 주어는?}
%     \begin{kfChoices}
%       \kfChoice 바람은
%       ...
%     \end{kfChoices}
%   \end{kfSentenceUnit}
% ============================================================

% kfSentenceBox 매크로는 passage-note.tex에 정의됨

\newenvironment{kfSentenceUnit}[2]%
  {\par\smallskip\needspace{6\baselineskip}%
   \noindent\begin{minipage}{\linewidth}%
   \samepage%
   \kfSentenceBox{문장 #1}{#2}%
   \par\smallskip%
  }%
  {\end{minipage}\par\smallskip}

% 같은 문장에 묶인 소문항 (문장 안내 + 발문)
\newcommand{\kfSentQ}[2]{%
  \par\noindent\textbf{#1}~#2\par\smallskip%
}

% 헤더 없이 문장만 있는 묶음 (sentence_ref가 없는 일반 문항)
\newenvironment{kfSentencelessUnit}%
  {\par\smallskip\noindent\begin{minipage}{\linewidth}\samepage}%
  {\end{minipage}\par\smallskip}

% ============================================================
% chapter-cover.tex (v6)
% 챕터 진입 페이지 (표지) — 하단에 영역 목차 추가
% 사용:
%   \kfChapterCover{1}{근대성과 자아}{Chapter 1}{이번 챕터에서는 ...}
%   \begin{kfChapterAreaList}
%     \kfChapterAreaItem{어휘}{kfAreaVocab}{핵심 어휘 12개}
%     ...
%   \end{kfChapterAreaList}
%
%   영역 목차가 없을 때는 \kfChapterCoverEnd 호출
% ============================================================

\newcommand{\kfChapterCover}[4]{%
  \clearpage%
  \thispagestyle{kfBlank}%
  \kfMarkChapter{Chapter #1. #2}%
  \kfChapterCoverBg%
  % v18.5: 챕터 표지 우상단에 로고
  \begin{tikzpicture}[remember picture, overlay]
    \node[anchor=north east, inner sep=0pt] at ($(current page.north east)+(-18mm,-18mm)$) {\kfLogoMd};
  \end{tikzpicture}%
  \vspace*{20mm}%
  \noindent
  \begin{tikzpicture}[remember picture, overlay]
    % 좌측 액센트 바
    \fill[kfPrimary] (current page.west) ++(12mm,25mm) rectangle ++(5mm, -50mm);
    % 챕터 번호 배지 (이중 원, 약간 작게)
    \node[anchor=north west, inner sep=0pt] at ($(current page.north west)+(24mm,-45mm)$) {%
      \begin{tikzpicture}
        \draw[kfPrimary, line width=1pt] (0,0) circle (10mm);
        \draw[kfPrimary, line width=0.4pt] (0,0) circle (8mm);
        \node[font=\normalsize\bfseries\color{kfPrimary}] at (0,2mm) {CHAPTER};
        \node[font=\LARGE\bfseries\color{kfPrimary}] at (0,-3mm) {#1};
      \end{tikzpicture}%
    };
  \end{tikzpicture}%
  \vspace{42mm}%
  \par\noindent\hspace{32mm}%
  \begin{minipage}{0.65\linewidth}%
    {\footnotesize\color{kfMute} #3}\\[4pt]
    {\LARGE\bfseries\color{kfPrimary} #2}\\[8pt]
    {\color{kfRule}\rule{50mm}{0.5pt}}\\[6pt]
    {\small\color{kfInk} #4}
  \end{minipage}%
  \par\vspace{14mm}%
}

% v6 / v18.5 / v18.6: 챕터 표지의 영역 목차 — 흰색 반투명 박스
% v18.6: 배경 가로선/도형과 안 겹치도록 TikZ overlay로 우하단 절대 위치
\makeatletter
% 영역 항목들을 모아 둘 박스
\newsavebox{\kf@arealistbox}
\newenvironment{kfChapterAreaList}{%
  \begin{lrbox}{\kf@arealistbox}%
  \begin{minipage}{0.62\linewidth}%
    \begin{tcolorbox}[%
      enhanced, colback=white, colframe=white,
      opacityback=0.92, opacityframe=0,
      boxrule=0pt, arc=4pt,
      width=\linewidth,
      top=8pt, bottom=8pt, left=10pt, right=10pt,
    ]%
      {\footnotesize\bfseries\color{kfMute} 이 챕터의 영역}\par\vspace{4pt}%
      {\color{kfRule}\hrule height 0.3pt}\par\vspace{4pt}%
      \begin{itemize}[leftmargin=14pt, itemsep=2pt, topsep=0pt, label={\color{kfPrimary!70}\textbullet}]%
}{%
      \end{itemize}%
    \end{tcolorbox}%
  \end{minipage}%
  \end{lrbox}%
  % v18.7: 배경 상단 가로선 ~ 하단 가로선 사이의 중간 영역에 배치
  % 가로선이 내용을 가리지 않도록 박스 중심이 두 선 사이 중점에 위치
  \begin{tikzpicture}[remember picture, overlay]%
    \node[anchor=center, inner sep=0pt]
      at ($(current page.south)+(0mm,90mm)$)
      {\usebox{\kf@arealistbox}};%
  \end{tikzpicture}%
  \vfill%
  \noindent\hfill\kfSeriesBadge\hspace{12mm}%
  \clearpage%
}
\makeatother

% 영역 항목: \kfChapterAreaItem{영역명}{kfArea색}{부제}
\makeatletter
\newcommand{\kfChapterAreaItem}[3]{%
  \item {\small\bfseries\color{#2} #1}%
  \def\kf@cci@sub{#3}%
  \ifx\kf@cci@sub\@empty\else
    {\small\color{kfMute}\, \textemdash\, #3}%
  \fi
}
\makeatother

% 영역 목록 없이 cover만 (legacy)
\newcommand{\kfChapterCoverEnd}{%
  \vfill%
  \noindent\hfill\kfSeriesBadge\hspace{12mm}%
  \clearpage%
}

% ============================================================
% area-header.tex (v16)
% 영역 시작 표제 — 시리즈별 차별화 라벨 + 큰 영역명 + 부제
% PAGE_BREAK_POLICY.md §1.4
%
% v16 (디자인 고도화 Phase 1):
%   - 시리즈별 라벨 박스 추가:
%     소쉬르 = AREA START (영역색 채움 박스)
%     프레게 = SECTION OPEN (영역색 테두리)
%     러셀   = RUSSELL (영역색 라이트 박스)
%     비트   = WITTGENSTEIN (회색 라이트 박스)
%   - 라벨 위치: 영역명 위 좌측
% ============================================================

\makeatletter

% 시리즈 라벨 텍스트
\newcommand{\kf@serieslabel}{%
  \ifcase\kf@series\relax
    AREA START%
  \or
    AREA START%
  \or
    SECTION OPEN%
  \or
    RUSSELL%
  \or
    WITTGENSTEIN%
  \fi
}

% 시리즈 라벨 박스 스타일 (#1 = 영역색)
\newcommand{\kf@labelbox}[1]{%
  \ifcase\kf@series\relax
    % 0/1 = 소쉬르: 영역색 채움
    \tikz[baseline=-2pt]{\node[fill=#1, text=white,
      rounded corners=3pt, inner xsep=6pt, inner ysep=2pt,
      font=\scriptsize\bfseries]{\kf@serieslabel};}%
  \or
    % 1 = 소쉬르: 영역색 채움
    \tikz[baseline=-2pt]{\node[fill=#1, text=white,
      rounded corners=3pt, inner xsep=6pt, inner ysep=2pt,
      font=\scriptsize\bfseries]{\kf@serieslabel};}%
  \or
    % 2 = 프레게: 영역색 테두리
    \tikz[baseline=-2pt]{\node[draw=#1, line width=0.6pt, text=#1,
      rounded corners=3pt, inner xsep=6pt, inner ysep=2pt,
      font=\scriptsize\bfseries]{\kf@serieslabel};}%
  \or
    % 3 = 러셀: 영역색 라이트 박스
    \tikz[baseline=-2pt]{\node[fill=#1!12, text=#1,
      rounded corners=2pt, inner xsep=5pt, inner ysep=1.5pt,
      font=\scriptsize\bfseries]{\kf@serieslabel};}%
  \or
    % 4 = 비트: 회색 미니 박스
    \tikz[baseline=-2pt]{\node[fill=kfMute!10, text=kfMute,
      rounded corners=2pt, inner xsep=5pt, inner ysep=1.5pt,
      font=\scriptsize\bfseries]{\kf@serieslabel};}%
  \fi
}

% v17: 시리즈+영역별 Canva 배경 자동 매핑
% \kf@hasbg=1로 set되면 \kfAreaStart에서 라벨 박스 출력 생략
\def\kf@areaVocab{어휘}\def\kf@areaGrammar{문법}\def\kf@areaConcept{개념}
\def\kf@areaLit{문학}\def\kf@areaNonlit{비문학}
\newif\ifkf@bgon
% 파일 있으면 배경 적용 + boolean true. 없으면 nothing.
\newcommand{\kfBgSet}[1]{%
  \IfFileExists{#1.pdf}{\kfPageBg{#1}\kf@bgontrue}{}%
}
\newcommand{\kf@trybg}[1]{%
  % v18.4: 영역 배경 PDF 비활성화 — 큰 도형이 본문 영역을 침범해
  % "영역 헤더 후 빈 페이지" 문제를 일으킴. 라벨 박스로만 영역 표시.
  \kf@bgonfalse%
}

% Note: 러셀(rus_*)/비트(wit_*) 배경 PDF 미존재 시 \kfPageBg가 에러 → \IfFileExists로 가드
% (현재는 파일이 모두 존재한다고 가정. 부분 제공 시 cls 수정 필요)

\newcommand{\kfAreaStart}[3]{%
  \clearpage%
  \thispagestyle{fancy}%
  \kfMarkArea{#1}%
  \kf@trybg{#1}%
  \vspace*{1mm}%
  % 배경 없으면 라텍스 라벨 박스 출력 (배경 있으면 일러스트가 라벨 역할)
  \ifkf@bgon
    \vspace*{4mm}%
  \else
    \noindent\kf@labelbox{#2}\par\vspace{4pt}%
  \fi
  % 영역명 + 부제
  \noindent
  \begin{tikzpicture}
    \node[anchor=west, inner sep=0pt] (bar) at (0,0) {\color{#2}\rule{3pt}{22pt}};
    \node[anchor=base west, inner sep=0pt, xshift=8pt, yshift=2pt] (lbl) at (bar.east |- 0,0) {%
      {\LARGE\bfseries\color{#2} #1}%
    };
    \def\kf@areasub{#3}%
    \ifx\kf@areasub\@empty\else
      \node[anchor=base west, inner sep=0pt, xshift=8pt, font=\normalsize\color{kfMute}]
        at (lbl.base east) {\,\textperiodcentered\, #3};%
    \fi
  \end{tikzpicture}%
  \par\vspace{2pt}
  {\color{#2!70}\hrule height 1pt}%
  \par\vspace{1pt}%
  {\color{kfRule}\hrule height 0.3pt}%
  \par\vspace{4pt}%
  \needspace{6\baselineskip}\nopagebreak[4]%
  \global\kf@areaJustOpenedtrue% v18.5: 첫 컨텐츠 needspace 무시
}
\makeatother

% 시리즈/영역 단축 매크로
\newcommand{\kfStartVocab}[1]{\kfAreaStart{어휘}{kfAreaVocab}{#1}}
\newcommand{\kfStartGrammar}[1]{\kfAreaStart{문법}{kfAreaGrammar}{#1}}
\newcommand{\kfStartConcept}[1]{\kfAreaStart{개념}{kfAreaConcept}{#1}}
\newcommand{\kfStartLit}[1]{\kfAreaStart{문학}{kfAreaLiterature}{#1}}
\newcommand{\kfStartNonlit}[1]{\kfAreaStart{비문학}{kfAreaNonlit}{#1}}
\newcommand{\kfStartWeekly}[1]{\kfAreaStart{실력 확인}{kfAreaWeekly}{#1}}

% ============================================================
% section-header.tex
% 섹션 시작 라벨 헤더 — [지문] / [활동] / [문제] 라벨
% 좌측 액센트 바 + 라벨 + 부제 (영역 액센트 색)
%
% 사용:
%   \kfSecHeader{kfAreaLiterature}{지문}{빼앗긴 들에도 봄은 오는가 / 이상화}
%   \kfSecHeader{kfAreaConcept}{활동}{내용 확인}
%   \kfSecHeader{kfAreaNonlit}{문제}{}
%
% 헤더 직후 페이지 분할 금지: \needspace{4\baselineskip} + \nopagebreak
% ============================================================

\providecommand{\kfSecHeader}[3]{%
  \par\needspace{10\baselineskip}\filbreak%
  \vspace{3pt}%
  \noindent
  \begin{tikzpicture}
    \node[anchor=west, inner sep=0pt] (bar) at (0,0) {\color{#1}\rule{3pt}{14pt}};
    \node[anchor=west, inner sep=0pt, xshift=6pt, font=\normalsize\bfseries\color{#1}]
       (lbl) at (bar.east) {[\,#2\,]};
    \node[anchor=west, inner sep=0pt, xshift=4pt, font=\small\color{kfMute}]
       at (lbl.east) {#3};
  \end{tikzpicture}%
  \par\vspace{1pt}%
  {\color{#1!50}\hrule height 0.4pt}%
  \par\vspace{2pt}\nopagebreak%
}

% 영역별 단축 매크로
\newcommand{\kfSecPassageVocab}[1]{\kfSecHeader{kfAreaVocab}{지문}{#1}}
\newcommand{\kfSecPassageGrammar}[1]{\kfSecHeader{kfAreaGrammar}{지문}{#1}}
\newcommand{\kfSecPassageConcept}[1]{\kfSecHeader{kfAreaConcept}{지문}{#1}}
\newcommand{\kfSecPassageLit}[1]{\kfSecHeader{kfAreaLiterature}{지문}{#1}}
\newcommand{\kfSecPassageNonlit}[1]{\kfSecHeader{kfAreaNonlit}{지문}{#1}}
\newcommand{\kfSecPassageWeekly}[1]{\kfSecHeader{kfAreaWeekly}{지문}{#1}}

\newcommand{\kfSecActVocab}[1]{\kfSecHeader{kfAreaVocab}{활동}{#1}}
\newcommand{\kfSecActGrammar}[1]{\kfSecHeader{kfAreaGrammar}{활동}{#1}}
\newcommand{\kfSecActConcept}[1]{\kfSecHeader{kfAreaConcept}{활동}{#1}}
\newcommand{\kfSecActLit}[1]{\kfSecHeader{kfAreaLiterature}{활동}{#1}}
\newcommand{\kfSecActNonlit}[1]{\kfSecHeader{kfAreaNonlit}{활동}{#1}}
\newcommand{\kfSecActWeekly}[1]{\kfSecHeader{kfAreaWeekly}{활동}{#1}}

\newcommand{\kfSecQVocab}[1]{\kfSecHeader{kfAreaVocab}{문제}{#1}}
\newcommand{\kfSecQGrammar}[1]{\kfSecHeader{kfAreaGrammar}{문제}{#1}}
\newcommand{\kfSecQConcept}[1]{\kfSecHeader{kfAreaConcept}{문제}{#1}}
\newcommand{\kfSecQLit}[1]{\kfSecHeader{kfAreaLiterature}{문제}{#1}}
\newcommand{\kfSecQNonlit}[1]{\kfSecHeader{kfAreaNonlit}{문제}{#1}}
\newcommand{\kfSecQWeekly}[1]{\kfSecHeader{kfAreaWeekly}{문제}{#1}}


\begin{document}

\thispagestyle{kfBlank}
\kfBookCoverBg
\vspace*{\kfBookCoverVspace}
\begin{center}
{\kfLogoLg}\\[14pt]
{\fontsize{72}{84}\selectfont\bfseries\kfBookCoverColor 소쉬르}\\[18pt]
{\fontsize{28}{32}\selectfont\kfBookCoverSubColor \textsc{Level} \textbf{3}}\\[22pt]
{\large\kfBookCoverSubColor 제 1 권 \quad\textbar\quad 챕터 1\textendash{} 4}
\end{center}
\vfill
\hfill\kfSeriesBadge\hspace{15mm}
\clearpage
\kfChapterCover{1}{소쉬르3 (챕터1)}{Chapter 1}{이번 챕터의 학습 내용을 시작합니다.}
\begin{kfChapterAreaList}
\kfChapterAreaItem{어휘}{kfAreaVocab}{나를 표현하는 말들}
\kfChapterAreaItem{개념}{kfAreaConcept}{배경지식 넓히기: 자아 성찰의 중요성}
\kfChapterAreaItem{문학}{kfAreaLiterature}{「거울 앞에 서면」}
\kfChapterAreaItem{비문학}{kfAreaNonlit}{나의 장점과 단점}
\kfChapterAreaItem{실력확인}{kfAreaWeekly}{}
\end{kfChapterAreaList}

\kfStartVocab{나를 표현하는 말들}


\begin{kfVocab}
\kfVocabItem{1}{자신의 마음이나 행동을 깊이 살펴보고 반성함}
\kfVocabItem{2}{마음속이나 정신의 안쪽 세계}
\kfVocabItem{3}{자기 자신을 깨닫는 의식의 주체}
\kfVocabItem{4}{개인이 가진 고유한 행동 양식이나 특성}
\kfVocabItem{5}{남보다 뛰어나거나 좋은 점}
\kfVocabItem{6}{부족하거나 모자라는 점}
\end{kfVocab}


\kfStartConcept{배경지식 넓히기: 자아 성찰의 중요성}


\kfPassageNoteNT{%
거울이 우리에게 알려주는 것

우리는 매일 거울을 봅니다. 거울 속에는 우리의 겉모습이 비춰집니다. 머리 모양이 어떤지, 옷이 잘 어울리는지 확인하죠. 하지만 거울이 보여주는 것은 우리의 외모뿐일까요?

옛날부터 많은 사람들은 거울을 '자아 성찰'의 도구로 사용했습니다. 자아 성찰이란 자신의 마음과 행동을 깊이 살펴보고 반성하는 것을 말합니다. 마치 거울이 우리의 겉모습을 비추듯이, 우리는 스스로를 돌아보며 내면의 모습을 살펴볼 수 있습니다.

자신을 알아가는 것은 매우 중요합니다. 내가 어떤 사람인지, 무엇을 잘하고 무엇이 부족한지 알아야 더 나은 사람으로 성장할 수 있기 때문입니다.

오늘부터 거울을 볼 때마다 겉모습뿐만 아니라 내면의 나도 함께 살펴보는 시간을 가져보는 것은 어떨까요?
\par\medskip\begin{itemize}[leftmargin=18pt, itemsep=2pt, topsep=2pt, label=\textbullet]
\item 겉모습 — 거울로 머리 모양·옷차림 확인
\item 내면(자아 성찰) — 오늘 한 말과 행동을 떠올리며 반성
\item 성장 — 부족한 점을 알면 더 나은 사람으로 변화
\end{itemize}
}

\kfPassageNoteNT{%
* 중심 문장: 문단에서 가장 중요하고 중심이 되는 문장입니다. 주로 문단의 처음이나 끝에 위치하지만, 글의 구성에 따라 위치가 달라질 수 있습니다. * 뒷받침 문장: 중심 문장을 더 쉽게 이해할 수 있도록 예시나 까닭을 들어 자세히 설명하는 문장입니다.

문단을 읽을 때 전체 내용을 포괄하는 문장을 찾으면 중심 문장을, 이를 구체적으로 설명하는 문장을 찾으면 뒷받침 문장을 쉽게 파악할 수 있습니다.
\par\medskip\begin{itemize}[leftmargin=18pt, itemsep=2pt, topsep=2pt, label=\textbullet]
\item 중심 문장 — '운동은 건강에 좋다'(문단 전체를 포괄)
\item 뒷받침(예시) — '예를 들어 매일 30분 걷기만 해도 심장이 튼튼해진다'
\item 뒷받침(까닭) — '왜냐하면 근육이 자극되어 혈액 순환이 좋아지기 때문이다'
\end{itemize}
}

\kfActivityBg \par\kfMaybeNeedspace{24\baselineskip}\noindent{\kfTipMark\,\,}{\small\color{kfInk} 다음 빈칸에 어울리는 어휘를 적어 보세요.}\par\nopagebreak[4]\smallskip\nopagebreak[4]
\begin{kfBlankList}
\kfBlankItem 나는 거울을 보며 나의 \kfFillBlank[42mm]{}을 깊이 들여다보았다.   

(마음속이나 정신의 안쪽을 가리키는 말)
\kfBlankItem 친구의 조언을 듣고 나의 행동을 \kfFillBlank[42mm]{}하는 시간을 가졌다.   

(자신을 깊이 살펴보고 반성하는 것)
\end{kfBlankList}

\kfActivityBg \par\kfMaybeNeedspace{24\baselineskip}\noindent{\kfTipMark\,\,}{\small\color{kfInk} 다음 설명이 맞으면 ○, 틀리면 ×를 표시해 보세요.}\par\nopagebreak[4]\smallskip\nopagebreak[4]
\begin{kfOXList}
\kfOXItem 거울은 우리의 겉모습만 보여준다. ( O / X ) \kfOX
\kfOXItem 자아 성찰은 자신을 깊이 살펴보고 반성하는 것이다. ( O / X ) \kfOX
\kfOXItem 자신의 장점과 단점을 아는 것은 성장에 도움이 된다. ( O / X ) \kfOX
\end{kfOXList}

\kfSecHeader{kfAreaConcept}{문제}{}

\begin{kfQBlock}{01}{객관식}
\kfQStem{\textbackslash{}<보기\textbackslash{}>에서 중심 문장은 몇 번 문장인가요?}
\begin{kfEvidence}① 좋은 친구는 우리 삶에 꼭 필요한 존재입니다.  \\\relax
 ② 힘들 때 위로해주고 격려해줍니다.  \\\relax
 ③ 기쁜 일이 있을 때 함께 기뻐해줍니다.  \\\relax
④ 또한 잘못된 길로 가려 할 때 바른 길을 알려줍니다.\end{kfEvidence}
\begin{kfChoices}
\kfChoice ① 번 문장
\kfChoice ② 번 문장
\kfChoice ③ 번 문장
\kfChoice ④ 번 문장
\end{kfChoices}
\end{kfQBlock}
\begin{kfQBlock}{02}{객관식}
\kfQStem{\textbackslash{}<보기\textbackslash{}>에서 중심 문장은 몇 번 문장인가요?}
\begin{kfEvidence}① 물은 우리 몸의 노폐물을 밖으로 내보내 줍니다.  \\\relax
 ② 또한 체온을 조절하는 중요한 역할도 합니다.  \\\relax
 ③ 피부를 촉촉하게 유지하는 데도 도움이 됩니다.  \\\relax
④ 따라서 우리는 물을 자주 마셔야 합니다.\end{kfEvidence}
\begin{kfChoices}
\kfChoice ① 번 문장
\kfChoice ② 번 문장
\kfChoice ③ 번 문장
\kfChoice ④ 번 문장
\end{kfChoices}
\end{kfQBlock}


\kfStartLit{「거울 앞에 서면」}


\kfPassageNote{}{%
거울 앞에 서면 \\[4pt]
또 다른 내가 있어요 \\[4pt]
웃으면 따라 웃고 \\[4pt]
찡그리면 따라 찡그려요 \\[14pt]
거울 속 친구는 \\[4pt]
항상 나와 함께해요 \\[4pt]
내가 행복할 때도 \\[4pt]
슬플 때도 곁에 있어요 \\[14pt]
가끔은 물어봐요 \\[4pt]
"너는 정말 나일까?" \\[4pt]
거울 속 나는 대답해요 \\[4pt]
"나는 네가 보는 너야" \\[14pt]
오늘도 거울을 봐요 \\[4pt]
나를 알아가기 위해 \\[4pt]
겉모습 너머에 있는 \\[4pt]
진짜 나를 찾아가요
}


\kfActivityBg \par\kfMaybeNeedspace{18\baselineskip}\noindent{\kfTipMark\,\,}{\small\color{kfInk} 주어진 안내에 맞춰 자유롭게 글로 표현해 보세요.}\par\nopagebreak[4]\smallskip\nopagebreak[4]
\kfWritingPrompt{'거울 속의 나'를 읽고, 여러분도 자신을 돌아보는 짧은 시를 써보세요. 아래의 형식을 참고하여 4행 정도로 작성해보세요.}
\kfWriting{답안 작성}{8}

\kfSecHeader{kfAreaLiterature}{문제}{}

\begin{kfQBlock}{01}{객관식}
\kfQStem{이 시에서 화자가 거울을 통해 만나는 대상은 무엇인가요?}
\begin{kfChoices}
\kfChoice 자신의 친구
\kfChoice 또 다른 자신
\kfChoice 상상의 인물
\kfChoice 가족 구성원
\end{kfChoices}
\end{kfQBlock}
\begin{kfQBlock}{02}{객관식}
\kfQStem{이 시의 전체적인 분위기로 가장 적절한 것은?}
\begin{kfChoices}
\kfChoice 슬프고 어둡고 우울한 분위기
\kfChoice 화나고 따져 묻는 분노의 분위기
\kfChoice 차분하고 성찰적인 분위기
\kfChoice 들뜨고 흥겹게 신나는 분위기
\end{kfChoices}
\end{kfQBlock}
\begin{kfQBlock}{03}{서술형}
\kfQStem{3연에서 화자가 거울 속 자신에게 던진 질문은 무엇인가요?}
\kfAnswerNote{답안}{4}
\end{kfQBlock}
\begin{kfQBlock}{04}{객관식}
\kfQStem{이 시에서 '거울'이 상징하는 의미로 가장 적절한 것은?}
\begin{kfChoices}
\kfChoice 단순한 생활용품
\kfChoice 자아 성찰의 도구
\kfChoice 아름다움의 기준
\kfChoice 시간의 흐름.
\end{kfChoices}
\end{kfQBlock}
\begin{kfQBlock}{05}{서술형}
\kfQStem{"나는 네가 보는 너야"라는 거울 속 '나'의 대답이 의미하는 바를 설명하시오.}
\kfAnswerNote{답안}{4}
\end{kfQBlock}
\begin{kfQBlock}{06}{서술형}
\kfQStem{마지막 연에서 화자가 찾고자 하는 '진짜 나'는 무엇을 의미하는지 자신의 생각을 써보세요.}
\kfAnswerNote{답안}{4}
\end{kfQBlock}


\kfStartNonlit{나의 장점과 단점}


\kfPassageNote{나의 장점과 단점}{%
사람은 누구나 장점과 단점을 가지고 있습니다. 장점은 남들보다 뛰어나거나 잘하는 점을 말하고, 단점은 부족하거나 고쳐야 할 점을 말합니다. 그런데 많은 사람들이 자신의 장점과 단점을 정확히 알지 못합니다. 왜 그럴까요?

첫째, 우리는 자신을 객관적으로 보기 어렵습니다. 자신의 행동이나 성격을 평가할 때 감정이 개입되기 쉽기 때문입니다. 예를 들어, 실수를 했을 때 '나는 원래 이런 사람이야'라고 지나치게 자책하거나, 반대로 '이번만 실수한 거야'라고 가볍게 넘기는 경우가 있습니다.

둘째, 다른 사람과 비교하는 습관 때문입니다. 친구가 수학을 잘한다고 해서 자신이 수학을 못한다고 생각하거나, 반대로 친구보다 그림을 잘 그린다고 해서 자신이 예술적 재능이 뛰어나다고 착각할 수 있습니다. 하지만 진정한 장점과 단점은 남과의 비교가 아닌, 자신의 다른 능력들과 비교해서 찾아야 합니다.

그렇다면 어떻게 자신의 장점과 단점을 찾을 수 있을까요? 먼저, 평소 자신이 좋아하고 즐겁게 하는 일이 무엇인지 생각해 보세요. 대부분의 사람들은 자신이 잘하는 일을 좋아합니다. 또한 주변 사람들의 이야기를 귀담아들어 보세요. 가족이나 친구들은 여러분이 보지 못하는 장점을 발견할 수 있습니다.

마지막으로, 자신의 단점을 부끄러워하지 마세요. 단점을 아는 것은 성장의 시작입니다. 오늘의 단점은 노력을 통해 내일의 장점이 될 수 있습니다. 중요한 것은 자신을 있는 그대로 받아들이고, 더 나은 사람이 되기 위해 꾸준히 노력하는 자세입니다.
}


\kfActivityBg \par\needspace{15\baselineskip}
\kfSecHeader{kfAreaNonlit}{활동}{문장 독해}
\par\kfMaybeNeedspace{32\baselineskip}\noindent{\kfTipMark\,\,}{\small\color{kfInk} 각 문장을 읽고 물음에 답하세요.}\par\nopagebreak[4]\smallskip\nopagebreak[4]
\begin{kfSentenceUnit}{1}{사람은 가지고 장점과 단점을}
\kfSentQ{1.}{(1) 이 문장에서 '누가'에 해당하는 것은 무엇인가요?}
\begin{kfChoices}
\kfChoice 사람은
\kfChoice 가지고
\kfChoice 장점과
\kfChoice 단점을
\end{kfChoices}
\kfSentQ{1.}{(2) 이 문장에서 '무엇을'에 해당하는 것은 무엇인가요?}
\begin{kfChoices}
\kfChoice 사람은
\kfChoice 누구나
\kfChoice 장점과 단점을
\kfChoice 가지고 있습니다
\end{kfChoices}
\kfSentQ{1.}{(3) 이 문장에서 '어찌하다'에 해당하는 것은 무엇인가요?}
\begin{kfChoices}
\kfChoice 사람은
\kfChoice 누구나
\kfChoice 장점과 단점을
\kfChoice 가지고 있습니다
\end{kfChoices}
\end{kfSentenceUnit}

\begin{kfSentenceUnit}{2}{많은 많은 사람들이 장점과 단점을}
\kfSentQ{2.}{(1) 이 문장에서 '누가'에 해당하는 것은 무엇인가요?}
\begin{kfChoices}
\kfChoice 많은
\kfChoice 많은 사람들이
\kfChoice 장점과
\kfChoice 단점을
\end{kfChoices}
\kfSentQ{2.}{(2) 이 문장에서 '무엇을'에 해당하는 것은 무엇인가요?}
\begin{kfChoices}
\kfChoice 자신의
\kfChoice 장점과
\kfChoice 정확히
\kfChoice 자신의 장점과 단점을
\end{kfChoices}
\kfSentQ{2.}{(3) 이 문장에서 '어찌하다'에 해당하는 것은 무엇인가요?}
\begin{kfChoices}
\kfChoice 정확히
\kfChoice 누구나
\kfChoice 단점을
\kfChoice 알지 못한다
\end{kfChoices}
\end{kfSentenceUnit}

\kfActivityBg \kfSecHeader{kfAreaNonlit}{활동}{지시어 연결}
\par\kfMaybeNeedspace{18\baselineskip}\noindent{\kfTipMark\,\,}{\small\color{kfInk} 다음 문장에서 지시어가 가리키는 대상을 찾으세요.}\par\nopagebreak[4]\smallskip\nopagebreak[4]

\kfActivityBg \par\kfMaybeNeedspace{24\baselineskip}\noindent{\kfTipMark\,\,}{\small\color{kfInk} 빈칸에 알맞은 말을 채워 보세요.}\par\nopagebreak[4]\smallskip\nopagebreak[4]
\begin{kfBlankList}
\kfBlankItem 요약문 완성하기: 각 문단의 중심 내용을 잘 읽고, 빈칸에 들어갈 알맞은 말을 지문에서 찾아 쓰세요.

1문단: 사람은 누구나 남들보다 뛰어난 \kfFillBlank[42mm]{}과/와 고쳐야 할 \kfFillBlank[42mm]{}을/를 가지고 있다. 2문단: 우리는 감정 때문에 스스로를 \kfFillBlank[42mm]{}으로 보기 어렵다. 3문단: 다른 사람과 \kfFillBlank[42mm]{}하는 습관은 나의 진짜 능력을 알기 어렵게 만든다. 4문단: 자신이 좋아하고 \kfFillBlank[42mm]{} 하는 일을 생각해보면 장점을 찾을 수 있다. 5문단: 단점을 \kfFillBlank[42mm]{}하지 말고 더 나은 사람이 되기 위해 노력해야 한다.
\end{kfBlankList}

\kfActivityBg \par\kfMaybeNeedspace{18\baselineskip}\noindent{\kfTipMark\,\,}{\small\color{kfInk} 주어진 안내에 맞춰 자유롭게 글로 표현해 보세요.}\par\nopagebreak[4]\smallskip\nopagebreak[4]
\kfWritingPrompt{지문에서 배운 방법을 활용하여 나의 장점 한 가지와 단점 한 가지를 찾아보고, 왜 그렇게 생각하는지 구체적인 이유와 함께 글로 써 보세요.}
\kfWriting{답안 작성}{8}

\kfSecHeader{kfAreaNonlit}{문제}{}

\begin{kfQBlock}{01}{객관식}
\kfQStem{윗글의 '그것'이 가리키는 대상으로 가장 알맞은 것은?}
\begin{kfChoices}
\kfChoice 사람들
\kfChoice 거울
\kfChoice 외모
\kfChoice 겉모습
\end{kfChoices}
\end{kfQBlock}
\begin{kfQBlock}{02}{객관식}
\kfQStem{윗글의 '이것'이 가리키는 대상으로 가장 알맞은 것은?}
\begin{kfChoices}
\kfChoice 사람들
\kfChoice 거울
\kfChoice 도구
\kfChoice 자아 성찰
\end{kfChoices}
\end{kfQBlock}
\begin{kfQBlock}{03}{객관식}
\kfQStem{이 글의 중심 주제로 가장 적절한 것은?}
\begin{kfChoices}
\kfChoice 다른 사람의 장단점을 정확히 평가하는 방법
\kfChoice 자신의 감정을 다스려 객관성을 키우는 방법
\kfChoice 자신의 장단점을 파악하는 방법
\kfChoice 좋아하는 일을 즐겁게 찾아 직업으로 삼는 방법
\end{kfChoices}
\end{kfQBlock}
\begin{kfQBlock}{04}{객관식}
\kfQStem{글쓴이가 제시한 '자신을 객관적으로 보기 어려운 이유'로 적절하지 \uline{\underline{않은}} 것은?}
\begin{kfChoices}
\kfChoice 감정이 개입되기 쉽기 때문
\kfChoice 다른 사람과 비교하기 때문
\kfChoice 거울을 자주 보지 않기 때문
\kfChoice 자책하거나 가볍게 넘기기 때문
\end{kfChoices}
\end{kfQBlock}
\begin{kfQBlock}{05}{객관식}
\kfQStem{글쓴이가 "오늘의 단점은 노력을 통해 내일의 장점이 될 수 있습니다"라는 말을 한 이유로 가장 알맞은 것은 무엇인가요?}
\begin{kfChoices}
\kfChoice 단점을 알려면 다른 사람과 자주 비교해야 함을 강조하기 위해
\kfChoice 자신의 단점을 부끄러워하지 말고 꾸준히 노력하면 단점이 장점으로 바뀔 수 있음을 알리기 위해
\kfChoice 단점을 객관적으로 보려면 감정을 완전히 없애야 함을 설명하기 위해
\kfChoice 좋아하고 즐겁게 하는 일이 있다면 단점이 자동으로 사라짐을 알리기 위해
\end{kfChoices}
\end{kfQBlock}
\begin{kfQBlock}{06}{서술형}
\kfQStem{글쓴이는 "진정한 장점과 단점은 남과의 비교가 \underline{아닌}, 자신의 다른 능력들과 비교해서 찾아야 합니다"라고 했습니다. 이 말의 의미를 구체적인 예를 들어 설명하시오.}
\kfAnswerNote{답안}{4}
\end{kfQBlock}


\kfStartWeekly{실력확인 영역 학습}


\noindent\textit{\small 제한시간: 30분 / 총 12문항}\par\medskip
\par\noindent\textbf{[4\textasciitilde{}5] 다음 글을 읽고 물음에 답하시오.}\par\medskip
\kfPassageNote{지문 4}{%
거울 속의 나 \\[4pt]
거울 앞에 서면 \\[4pt]
또 다른 내가 있어요 \\[4pt]
웃으면 따라 웃고 \\[4pt]
찡그리면 따라 찡그려요 \\[4pt]
거울 속 친구는 \\[4pt]
항상 나와 함께해요 \\[4pt]
내가 행복할 때도 \\[4pt]
슬플 때도 곁에 있어요 \\[4pt]
가끔은 물어봐요 \\[4pt]
"너는 정말 나일까?" \\[4pt]
거울 속 나는 대답해요" \\[4pt]
나는 네가 보는 너야" \\[4pt]
오늘도 거울을 봐요 \\[4pt]
나를 알아가기 위해 \\[4pt]
겉모습 너머에 있는 \\[4pt]
진짜 나를 찾아가요
}
\par\noindent\textbf{[6\textasciitilde{}7] 다음 글을 읽고 물음에 답하시오.}\par\medskip
\kfPassageNote{지문 6}{%
나의 장점과 단점 사람은 누구나 장점과 단점을 가지고 있습니다. 장점은 남들보다 뛰어나거나 잘하는 점을 말하고, 단점은 부족하거나 고쳐야 할 점을 말합니다. 그런데 많은 사람들이 자신의 장점과 단점을 정확히 알지 못합니다. 왜 그럴까요? 첫째, 우리는 자신을 객관적으로 보기 어렵습니다. 자신의 행동이나 성격을 평가할 때 감정이 개입되기 쉽기 때문입니다. 예를 들어, 실수를 했을 때 '나는 원래 이런 사람이야'라고 지나치게 자책하거나, 반대로 '이번만 실수한 거야'라고 가볍게 넘기는 경우가 있습니다. 둘째, 다른 사람과 비교하는 습관 때문입니다. 친구가 수학을 잘한다고 해서 자신이 수학을 못한다고 생각하거나, 반대로 친구보다 그림을 잘 그린다고 해서 자신이 예술적 재능이 뛰어나다고 착각할 수 있습니다. 하지만 진정한 장점과 단점은 남과의 비교가 아닌, 자신의 다른 능력들과 비교해서 찾아야 합니다. 그렇다면 어떻게 자신의 장점과 단점을 찾을 수 있을까요? 먼저, 평소 자신이 좋아하고 즐겁게 하는 일이 무엇인지 생각해 보세요. 대부분 사람들은 자신이 잘하는 일을 좋아합니다. 또한 주변 사람들의 이야기를 귀담아들어 보세요. 가족이나 친구들은 여러분이 보지 못하는 장점을 발견할 수 있습니다. 마지막으로, 자신의 단점을 부끄러워하지 마세요. 단점을 아는 것은 성장의 시작입니다. 오늘의 단점은 노력을 통해 내일의 장점이 될 수 있습니다. 중요한 것은 자신을 있는 그대로 받아들이고, 더 나은 사람이 되기 위해 꾸준히 노력하는 자세입니다.
}
\begin{kfQBlock}{01}{객관식}
\kfQStem{다음 중 '성찰'의 의미로 가장 적절한 것은?}
\begin{kfChoices}
\kfChoice 겉모습을 꾸미는 것
\kfChoice 자신을 깊이 살펴보고 반성하는 것
\kfChoice 다른 사람을 관찰하는 것
\kfChoice 미래를 계획하는 것
\end{kfChoices}
\end{kfQBlock}

\begin{kfQBlock}{02}{객관식}
\kfQStem{다음 문장의 빈칸에 들어갈 알맞은 말은 무엇인가요?}
\begin{kfEvidence}"거울은 우리의 겉모습을 비추지만, 자아 성찰은 마음속 깊은 곳인 \kfFillBlank[30mm]{}을/를 들여다보게 합니다."\end{kfEvidence}
\begin{kfChoices}
\kfChoice 외면
\kfChoice 내면
\kfChoice 앞면
\kfChoice 뒷면
\end{kfChoices}
\end{kfQBlock}

\begin{kfQBlock}{03}{객관식}
\kfQStem{자아 성찰에 대한 설명으로 옳은 것은 무엇인가요?}
\begin{kfChoices}
\kfChoice 거울을 보며 외모를 꾸미는 것이다.
\kfChoice 자신의 잘못은 덮어두고 잘한 점만 생각하는 것이다.
\kfChoice 자신의 마음과 행동을 깊이 살피며 반성하는 것이다.
\kfChoice 다른 사람의 행동을 관찰하고 평가하는 것이다.
\end{kfChoices}
\end{kfQBlock}

\begin{kfQBlock}{04}{객관식}
\kfQStem{이 시에서 화자가 거울을 보는 이유는 무엇인가요?}
\begin{kfChoices}
\kfChoice 외모를 확인하기 위해
\kfChoice 진짜 자신을 찾아가기 위해
\kfChoice 친구를 만나기 위해
\kfChoice 시간을 보내기 위해
\end{kfChoices}
\end{kfQBlock}

\begin{kfQBlock}{05}{객관식}
\kfQStem{'거울'이 상징하는 의미로 가장 알맞은 것은 무엇인가요?}
\begin{kfChoices}
\kfChoice 시간을 알려주는 도구
\kfChoice 자아 성찰의 도구
\kfChoice 미래를 보여주는 도구
\kfChoice 친구를 사귀는 도구
\end{kfChoices}
\end{kfQBlock}

\begin{kfQBlock}{06}{객관식}
\kfQStem{위 글에서 자신의 장단점을 찾는 방법으로 제시되지 \uline{\underline{않은}} 것은?}
\begin{kfChoices}
\kfChoice 자신이 좋아하는 일 생각해보기
\kfChoice 주변 사람들의 이야기 듣기
\kfChoice 전문가의 검사 받기
\kfChoice 자신의 능력들끼리 비교하기
\end{kfChoices}
\end{kfQBlock}

\begin{kfQBlock}{07}{객관식}
\kfQStem{글쓴이가 말하는 '단점을 대하는 태도'로 가장 알맞은 것은 무엇인가요?}
\begin{kfChoices}
\kfChoice 남들에게 들키지 않도록 꽁꽁 숨긴다.
\kfChoice 단점은 고칠 수 없으니 포기한다.
\kfChoice 단점을 부끄러워하지 않고 노력해서 고치려고 한다.
\kfChoice 장점만 생각하고 단점은 아예 무시한다.
\end{kfChoices}
\end{kfQBlock}

\begin{kfQBlock}{08}{객관식}
\kfQStem{다음 문장에서 중심 문장과 뒷받침 문장의 관계를 올바르게 설명한 것은?}
\begin{kfEvidence}㉠ 독서는 우리에게 꼭 필요합니다. \\\relax
㉡ 독서는 상상력을 키워 줍니다. \\\relax
㉢ 독서는 다양한 지식을 얻게 해 줍니다.\end{kfEvidence}
\begin{kfChoices}
\kfChoice ㉠이 중심 문장, ㉡㉢이 뒷받침 문장
\kfChoice ㉡이 중심 문장, ㉠㉢이 뒷받침 문장
\kfChoice ㉢이 중심 문장, ㉠㉡이 뒷받침 문장
\kfChoice 모두 동등한 중심 문장
\end{kfChoices}
\end{kfQBlock}

\begin{kfQBlock}{09}{객관식}
\kfQStem{다음 빈칸에 들어갈 말로 가장 적절한 것은?}
\begin{kfEvidence}거울은 우리의 \kfFillBlank[42mm]{}을/를 비추지만, 자아 성찰은 우리의 \kfFillBlank[42mm]{}을/를 들여다보게 한다.\end{kfEvidence}
\begin{kfChoices}
\kfChoice 외모 \textbackslash{}- 내면
\kfChoice 과거 \textbackslash{}- 미래
\kfChoice 장점 \textbackslash{}- 단점
\kfChoice 행동 \textbackslash{}- 말
\end{kfChoices}
\end{kfQBlock}

\begin{kfQBlock}{10}{서술형}
\kfQStem{다음 문장에서 주어와 서술어의 호응이 잘못된 곳을 찾아 바르게 고치시오.}
\begin{kfEvidence}"거울을 보는 이유는 내 마음을 살핀다."\end{kfEvidence}
\kfAnswerNote{답안}{4}
\end{kfQBlock}

\begin{kfQBlock}{11}{서술형}
\kfQStem{이번 주에 배운 내용을 바탕으로, '나를 다시 보기'의 중요성에 대해 3문장으로 서술하시오.}
\kfAnswerNote{답안}{4}
\end{kfQBlock}




\clearpage

\kfChapterCover{2}{소쉬르3 (챕터2)}{Chapter 2}{이번 챕터의 학습 내용을 시작합니다.}
\begin{kfChapterAreaList}
\kfChapterAreaItem{어휘}{kfAreaVocab}{마음을 표현하는 말들}
\kfChapterAreaItem{개념}{kfAreaConcept}{배경지식 넓히기: 모든 일에는 이유가 있어요}
\kfChapterAreaItem{문학}{kfAreaLiterature}{「옛날 어느 숲속에 힘이 세지만 성격이 아…」}
\kfChapterAreaItem{비문학}{kfAreaNonlit}{감정이 우리에게 보내는 신호}
\end{kfChapterAreaList}

\kfStartVocab{마음을 표현하는 말들}


\begin{kfVocab}
\kfVocabItem{1}{어떤 일이나 현상에 대하여 일어나는 마음이나 기분}
\kfVocabItem{2}{어떤 일이 일어나게 된 까닭이나 근본이 된 일}
\kfVocabItem{3}{어떤 원인으로 결말이 생김. 또는 그런 결말}
\kfVocabItem{4}{균형이 맞게 바로잡거나 상황에 알맞게 맞춤}
\kfVocabItem{5}{이전에 자신이 한 일을 잘못된 것으로 느껴 뉘우침}
\kfVocabItem{6}{어려운 일이나 좋지 않은 상태를 풀어서 없애 버림}
\end{kfVocab}


\kfStartConcept{배경지식 넓히기: 모든 일에는 이유가 있어요}


\kfPassageNoteNT{%
원인과 결과의 짝꿍

세상에 일어나는 모든 일에는 이유가 있습니다. 배가 고픈 이유는 밥을 먹지 않았기 때문이고, 식물이 잘 자라는 이유는 물과 햇빛을 충분히 받았기 때문입니다. 이렇게 어떤 일이 일어나게 된 까닭을 '원인'이라고 하고, 그 원인 때문에 생겨난 일을 '결과'라고 합니다.

우리의 감정도 마찬가지입니다. 감정은 대개 어떤 이유와 함께 생깁니다. 예를 들어 친구에게 칭찬을 들어서 기분이 좋다면, 칭찬을 들은 일은 원인이 되고 기쁜 마음은 결과가 됩니다. 또 아끼던 장난감이 망가져서 속상하다면, 장난감이 망가진 것이 원인이고 속상한 감정은 결과가 되는 것입니다.

나의 감정을 잘 이해하려면 원인과 결과가 서로 짝을 이루는 '인과 관계'를 생각해보는 습관이 필요합니다. "나 지금 왜 화가 났지?" 하고 스스로에게 물어보면, 화가 난 이유인 원인을 찾을 수 있습니다. 원인을 알면 그 마음을 어떻게 조절해야 할지도 알 수 있게 됩니다. 문제를 해소하는 방법을 찾기 쉬워지기 때문입니다.
\par\medskip\begin{itemize}[leftmargin=18pt, itemsep=2pt, topsep=2pt, label=\textbullet]
\item 원인 → 결과(자연) — 햇빛·물(원인) → 식물 잘 자람(결과)
\item 원인 → 결과(감정) — 칭찬을 들음(원인) → 기쁨(결과)
\item 감정 조절 — '왜 화가 났지?'로 원인을 찾으면 해결책도 보임
\end{itemize}
}

\kfPassageNoteNT{%
원인은 어떤 일이 일어나게 된 까닭이고, 결과는 그 때문에 일어난 일입니다. 원인과 결과는 짝꿍처럼 붙어 다닙니다.
}

\kfActivityBg \par\kfMaybeNeedspace{24\baselineskip}\noindent{\kfTipMark\,\,}{\small\color{kfInk} 다음 빈칸에 어울리는 어휘를 적어 보세요.}\par\nopagebreak[4]\smallskip\nopagebreak[4]
\begin{kfBlankList}
\kfBlankItem 배가 아픈 \kfFillBlank[42mm]{}을 찾아보니, 찬 음식을 너무 많이 먹어서였다.   (어떤 일이 일어나게 된 까닭)
\kfBlankItem 동생에게 화를 낸 뒤에, 나는 곧 나의 행동을 \kfFillBlank[42mm]{}했다.   (자신의 일을 잘못된 것으로 느껴 뉘우침)
\end{kfBlankList}

\kfActivityBg \par\kfMaybeNeedspace{24\baselineskip}\noindent{\kfTipMark\,\,}{\small\color{kfInk} 다음 설명이 맞으면 ○, 틀리면 ×를 표시해 보세요.}\par\nopagebreak[4]\smallskip\nopagebreak[4]
\begin{kfOXList}
\kfOXItem '원인'은 어떤 일이 일어난 결과를 말한다. ( O / X ) \kfOX
\kfOXItem 우리의 감정에는 생겨나는 이유가 있다. ( O / X ) \kfOX
\kfOXItem 감정의 원인을 알면 마음을 조절하는 데 도움이 된다. ( O / X ) \kfOX
\end{kfOXList}

\kfSecHeader{kfAreaConcept}{문제}{}

\begin{kfQBlock}{01}{객관식}
\kfQStem{민수가 오늘 아침에 늦잠을 잔 원인은 무엇인가요?}
\begin{kfChoices}
\kfChoice 가기 싫은 학교에 가야 해서
\kfChoice 밤늦게까지 텔레비전을 보아서
\kfChoice 머리맡 알람 시계가 고장 나서
\kfChoice 새벽부터 배가 너무 고파서
\end{kfChoices}
\end{kfQBlock}
\begin{kfQBlock}{02}{서술형}
\kfQStem{민수가 학교에 지각을 한 결과, 어떤 일이 일어났나요? ( 선생님께 \_\_\_\_\_\_\_\_\_\_\_\_\_\_\_\_\_\_\_\_\_\_\_\_\_\_\_\_\_\_\_\_\_\_\_\_\_\_\_ )}
\kfAnswerNote{답안}{4}
\end{kfQBlock}


\kfStartLit{「옛날 어느 숲속에 힘이 세지만 성격이 아…」}


\kfPassageNote{}{%
옛날 어느 숲속에 힘이 세지만 성격이 아주 급한 사자가 살았습니다. 사자는 조금만 기분이 나빠도 크게 소리를 지르며 화를 내곤 했습니다. 동물들은 언제 화를 낼지 모르는 사자를 피해 슬금슬금 도망을 다녔습니다.

어느 무더운 여름날 오후였습니다. 사자가 낮잠을 자고 있는데 어디선가 윙윙거리는 소리가 들렸습니다. 조그만 모기 한 마리가 사자의 콧등 주변을 날아다니고 있었습니다.

"감히 사자님의 낮잠을 방해하다니!"

사자는 몹시 화가 나서 앞발을 크게 휘둘렀습니다. 하지만 모기는 잽싸게 피했고, 사자의 날카로운 발톱은 자신의 코만 할퀴고 말았습니다. 화가 머리끝까지 난 사자는 더욱 거칠게 날뛰며 모기를 잡으려 했습니다. 으르렁거리는 소리에 숲속 동물들은 모두 숨을 죽이고 바들바들 떨었습니다.

한참을 날뛰던 사자는 숨이 차서 연못가에 주저앉았습니다. 그러다 물에 비친 자신의 얼굴을 보고 깜짝 놀랐습니다. 콧등과 볼은 자신이 휘두른 발톱에 긁혀 상처투성이였고, 화난 표정은 괴물처럼 무서워 보였습니다. 정작 모기는 윙윙거리며 유유히 날아가 버린 뒤였습니다.

"고작 모기 한 마리 때문에 내 얼굴을 이렇게 만들다니."

사자는 쓰라린 얼굴을 어루만지며 깊이 후회했습니다. 그리고 참을성 없이 터뜨린 화가 결국 자신을 아프게 한다는 사실을 깨달았습니다. 그날 이후 사자는 화가 나는 일이 생기면 마음속으로 숫자를 세며 기분을 가라앉히는 연습을 시작했습니다.
}


\kfActivityBg \par\kfMaybeNeedspace{18\baselineskip}\noindent{\kfTipMark\,\,}{\small\color{kfInk} 주어진 안내에 맞춰 자유롭게 글로 표현해 보세요.}\par\nopagebreak[4]\smallskip\nopagebreak[4]
\kfWritingPrompt{사자는 화가 날 때 '숫자 세기'를 하기로 했습니다. 여러분은 화가 나거나 속상할 때 어떻게 기분을 풀나요? 나만의 방법을 소개해 주세요.}
\kfWriting{답안 작성}{8}

\kfSecHeader{kfAreaLiterature}{문제}{}

\begin{kfQBlock}{01}{객관식}
\kfQStem{사자가 낮잠을 자다가 화가 난 원인은 무엇인가요?}
\begin{kfChoices}
\kfChoice 다른 사자가 공격해서
\kfChoice 배가 너무 고파서
\kfChoice 모기가 콧등 주변을 날아다녀서
\kfChoice 숲속 동물들이 시끄럽게 떠들어서
\end{kfChoices}
\end{kfQBlock}
\begin{kfQBlock}{02}{객관식}
\kfQStem{사자가 연못에 비친 자신의 모습을 보고 깨달은 점으로 가장 적절한 것은 무엇인가요?}
\begin{kfChoices}
\kfChoice 모기를 잡으려면 더 빠르게 발톱을 휘둘러야 한다.
\kfChoice 참을성 없이 터뜨린 화가 결국 자신의 얼굴을 아프게 한다.
\kfChoice 동물들이 무서워한 진짜 까닭은 내 큰 목소리 때문이었다.
\kfChoice 모기가 또 와도 무서운 표정을 지어 두면 다시는 다가오지 않는다.
\end{kfChoices}
\end{kfQBlock}
\begin{kfQBlock}{03}{서술형}
\kfQStem{사자가 화를 참지 못하고 날뛴 결과로 어떤 일이 일어났나요?}
\kfAnswerNote{답안}{4}
\end{kfQBlock}
\begin{kfQBlock}{04}{객관식}
\kfQStem{이야기의 마지막 부분에서 사자는 화가 날 때 어떻게 하기로 했나요?}
\begin{kfChoices}
\kfChoice 더 크게 소리를 지르기로 했다.
\kfChoice 모기를 끝까지 쫓아가기로 했다.
\kfChoice 동굴 속에 들어가서 숨기로 했다.
\kfChoice 마음속으로 숫자를 세며 기분을 가라앉히기로 했다.
\end{kfChoices}
\end{kfQBlock}
\begin{kfQBlock}{05}{서술형}
\kfQStem{만약 사자가 모기를 처음 봤을 때 화를 내지 않고 꼬리로 살짝 쫓아버렸다면 결과는 어떻게 달라졌을까요? 상상해서 적어보세요.}
\kfAnswerNote{답안}{4}
\end{kfQBlock}


\kfStartNonlit{감정이 우리에게 보내는 신호}


\kfPassageNote{감정이 우리에게 보내는 신호}{%
우리는 하루에도 수십 번씩 기분이 변합니다. 맛있는 밥을 먹으면 행복하고, 숙제가 많으면 짜증이 나기도 합니다. 어떤 사람들은 화를 내거나 슬퍼하는 것을 나쁜 것이라고 생각해서 꾹 참기도 합니다. 하지만 모든 감정은 우리 몸과 마음이 보내는 중요한 신호입니다. 마치 자동차가 다니는 도로의 신호등처럼, 감정도 우리가 어떻게 행동해야 할지 알려주는 역할을 합니다.

그렇다면 각각의 감정은 어떤 신호를 보내는 것일까요? 먼저 '화'는 나를 지켜달라는 신호입니다. 누군가 나를 부당하게 대하거나 내 물건을 함부로 가져갔을 때 우리는 화가 납니다. 이것은 나의 소중한 것을 보호해야 한다는 뜻입니다.

'슬픔'은 위로가 필요하다는 신호입니다. 아끼던 것을 잃어버리거나 사랑하는 사람과 헤어졌을 때 슬픔을 느낍니다. 이때 흘리는 눈물은 마음의 아픔을 씻어 내고, 주변 사람들에게 "나 지금 힘드니까 도와주세요"라고 말하는 것과 같습니다.

'두려움'이나 무서움은 위험을 피하라는 신호입니다. 캄캄한 골목길이나 사나운 개를 만났을 때 가슴이 두근거리는 것은 위험할 수 있으니 조심하거나 도망치라고 우리 몸이 경고를 보내는 것입니다.

이처럼 쓸모없는 감정은 하나도 없습니다. 감정이 보내는 신호를 무시하지 않고 잘 귀담아듣는다면, 우리는 어떤 상황에서도 더 지혜롭게 행동할 수 있을 것입니다.
}


\kfActivityBg \par\needspace{15\baselineskip}
\kfSecHeader{kfAreaNonlit}{활동}{문장 독해}
\par\kfMaybeNeedspace{32\baselineskip}\noindent{\kfTipMark\,\,}{\small\color{kfInk} 각 문장을 읽고 물음에 답하세요.}\par\nopagebreak[4]\smallskip\nopagebreak[4]
\begin{kfSentenceUnit}{1}{친구가 칭찬을 해 주어서 기분이 날아갈 듯 좋다}
\kfSentQ{1.}{이 문장에서 '원인'에 해당하는 것은 무엇인가요?}
\begin{kfChoices}
\kfChoice 친구가 칭찬을 해 주어서
\kfChoice 기분이 날아갈 듯 좋다
\end{kfChoices}
\end{kfSentenceUnit}

\begin{kfSentenceUnit}{2}{아끼던 장난감이 망가져서 눈물이 났다}
\kfSentQ{2.}{이 문장에서 '결과'에 해당하는 것은 무엇인가요?}
\begin{kfChoices}
\kfChoice 아끼던 장난감이 망가져서
\kfChoice 눈물이 났다
\end{kfChoices}
\end{kfSentenceUnit}

\begin{kfSentenceUnit}{3}{비가 많이 내려서 운동장에서 놀지 못했다}
\kfSentQ{3.}{이 문장에서 '원인'에 해당하는 것은 무엇인가요?}
\begin{kfChoices}
\kfChoice 비가 많이 내려서
\kfChoice 운동장에서 놀지 못했다
\end{kfChoices}
\end{kfSentenceUnit}

\kfActivityBg \kfSecHeader{kfAreaNonlit}{활동}{지시어 연결}
\par\kfMaybeNeedspace{18\baselineskip}\noindent{\kfTipMark\,\,}{\small\color{kfInk} 다음 문장에서 지시어가 가리키는 대상을 찾으세요.}\par\nopagebreak[4]\smallskip\nopagebreak[4]

\kfActivityBg \par\kfMaybeNeedspace{24\baselineskip}\noindent{\kfTipMark\,\,}{\small\color{kfInk} 빈칸에 알맞은 말을 채워 보세요.}\par\nopagebreak[4]\smallskip\nopagebreak[4]
\begin{kfBlankList}
\kfBlankItem 요약문 완성하기: 각 문단의 중심 내용을 잘 읽고, 빈칸에 들어갈 알맞은 말을 지문에서 찾아 쓰세요.

1문단: 감정은 우리 몸과 마음이 보내는 중요한 \kfFillBlank[42mm]{}이다. 2문단: '화'는 나의 소중한 것을 \kfFillBlank[42mm]{}해야 한다는 뜻이다. 3문단: '슬픔'은 \kfFillBlank[42mm]{}가 필요하다는 신호이다. 4문단: '두려움'은 \kfFillBlank[42mm]{}을 피하라고 우리 몸이 보내는 경고이다. 5문단: \kfFillBlank[42mm]{} 감정은 하나도 없으니, 감정의 소리를 잘 들어야 한다.
\end{kfBlankList}

\kfActivityBg \par\kfMaybeNeedspace{18\baselineskip}\noindent{\kfTipMark\,\,}{\small\color{kfInk} 주어진 안내에 맞춰 자유롭게 글로 표현해 보세요.}\par\nopagebreak[4]\smallskip\nopagebreak[4]
\kfWritingPrompt{최근에 느꼈던 강한 감정 하나를 골라, 그 감정이 나에게 어떤 신호를 보냈는지 적어보세요.}
\kfWriting{답안 작성}{8}

\kfActivityBg \par\kfMaybeNeedspace{18\baselineskip}\noindent{\kfTipMark\,\,}{\small\color{kfInk} 다음 활동을 풀어 보세요.}\par\nopagebreak[4]\smallskip\nopagebreak[4]
\begin{enumerate}[leftmargin=22pt, itemsep=6pt, topsep=4pt, label=\textbf{\arabic*.}]
\item 각 감정이 보내는 신호를 바르게 연결해 보세요.
\end{enumerate}

\kfSecHeader{kfAreaNonlit}{문제}{}

\begin{kfQBlock}{01}{객관식}
\kfQStem{윗글의 '이것'이 가리키는 대상으로 가장 알맞은 것은?}
\begin{kfChoices}
\kfChoice 물건
\kfChoice 우리
\kfChoice 화가 나는 것
\kfChoice 소중한 것
\end{kfChoices}
\end{kfQBlock}
\begin{kfQBlock}{02}{객관식}
\kfQStem{이 글의 중심 내용으로 가장 적절한 것은 무엇인가요?}
\begin{kfChoices}
\kfChoice 신호등이 어떻게 자동차를 멈추고 가게 하는지 그 원리를 알려 주는 글이다.
\kfChoice 화·슬픔·두려움 같은 감정이 우리에게 보내는 신호와 그 역할을 설명한 글이다.
\kfChoice 캄캄한 골목길과 사나운 개를 만났을 때 도망치는 방법을 알려 주는 글이다.
\kfChoice 슬플 때 흘리는 눈물의 성분을 과학적으로 분석한 글이다.
\end{kfChoices}
\end{kfQBlock}
\begin{kfQBlock}{03}{객관식}
\kfQStem{글쓴이는 감정을 무엇에 비유하여 설명했나요?}
\begin{kfChoices}
\kfChoice 맛있는 밥
\kfChoice 캄캄한 골목길
\kfChoice 도로의 신호등
\kfChoice 사나운 개
\end{kfChoices}
\end{kfQBlock}
\begin{kfQBlock}{04}{서술형}
\kfQStem{글쓴이는 "쓸모\underline{없는} 감정은 하나도 없습니다"라고 말했습니다. 내가 느끼기 싫은 '짜증'이나 '부끄러움' 같은 감정도 필요한 이유가 무엇일지 생각해서 적어보세요.}
\kfAnswerNote{답안}{4}
\end{kfQBlock}
\begin{kfQBlock}{05}{서술형}
\kfQStem{만약 우리에게 '두려움'이라는 감정이 아예 사라진다면 어떤 일이 벌어질까요? 인과 관계(원인과 결과)를 생각하며 적어보세요. * 원인: 두려움을 전혀 느끼지 못한다면 * 결과:}
\kfAnswerNote{답안}{4}
\end{kfQBlock}



\clearpage

\kfChapterCover{3}{소쉬르3 (챕터3)}{Chapter 3}{이번 챕터의 학습 내용을 시작합니다.}
\begin{kfChapterAreaList}
\kfChapterAreaItem{어휘}{kfAreaVocab}{대화를 돕는 말들}
\kfChapterAreaItem{문법}{kfAreaGrammar}{문장의 짜임}
\kfChapterAreaItem{개념}{kfAreaConcept}{배경지식 넓히기: 마음을 여는 열쇠, 공감}
\kfChapterAreaItem{문학}{kfAreaLiterature}{「숲속 마을에 깃털이 알록달록한 앵무새 '…」}
\kfChapterAreaItem{비문학}{kfAreaNonlit}{공감하며 대화하는 방법}
\kfChapterAreaItem{실력확인}{kfAreaWeekly}{}
\end{kfChapterAreaList}

\kfStartVocab{대화를 돕는 말들}


\begin{kfVocab}
\kfVocabItem{1}{마주 대하여 이야기를 주고받음}
\kfVocabItem{2}{귀를 기울여 남의 말을 주의 깊게 들음}
\kfVocabItem{3}{남의 감정, 의견, 주장 따위에 대하여 자기도 그렇다고 느낌}
\kfVocabItem{4}{그릇되게 해석하거나 뜻을 잘못 앎}
\kfVocabItem{5}{거짓이 없는 참된 마음}
\kfVocabItem{6}{어떤 자극에 대하여 일어나는 동작이나 태도}
\end{kfVocab}


\kfStartGrammar{문장의 짜임}


\kfPassageNoteNT{%
우리말 문장은 크게 세 가지 짜임으로 이루어져 있습니다. 문장의 뼈대를 잘 알면 글을 더 잘 이해할 수 있습니다.

* 무엇이 무엇이다: 주어의 정체나 이름을 설명합니다.        예) 나는 초등학생이다.   * 무엇이 어찌하다: 주어의 움직임이나 행동을 설명합니다.        예) 파도가 밀려온다.   * 무엇이 어떠하다: 주어의 성질이나 상태를 설명합니다.        예) 날씨가 매우 덥다.
\par\medskip\begin{itemize}[leftmargin=18pt, itemsep=2pt, topsep=2pt, label=\textbullet]
\item 무엇이 무엇이다: '나는 초등학생이다.'
\item 무엇이 어찌하다: '파도가 밀려온다.'
\item 무엇이 어떠하다: '날씨가 매우 덥다.'
\end{itemize}
}

\kfActivityBg \par\kfMaybeNeedspace{24\baselineskip}\noindent{\kfTipMark\,\,}{\small\color{kfInk} 다음 문장이 어떤 짜임인지 찾아 선으로 연결해 보세요.}\par\nopagebreak[4]\smallskip\nopagebreak[4]
\begin{kfConnect}
\kfPair{"부엉이는 새이다."}{⦁ 무엇이 무엇이다}
\kfPair{"친구가 소리 내어 웃는다."}{⦁ 무엇이 어찌하다}
\kfPair{"오늘 하늘이 정말 파랗다."}{⦁ 무엇이 어떠하다}
\end{kfConnect}


\kfStartConcept{배경지식 넓히기: 마음을 여는 열쇠, 공감}


\kfPassageNoteNT{%
귀는 두 개, 입은 한 개

우리의 귀는 두 개이고 입은 하나입니다. 옛 어른들은 이것을 두고 말하기보다는 듣기를 두 배 더 많이 하라는 뜻이라고 했습니다. 대화에서 가장 중요한 것은 내 말을 잘하는 것이 아니라 상대방의 말을 잘 들어주는 것이기 때문입니다.

상대방의 말을 잘 듣고 마음을 이해해 주는 것을 '공감'이라고 합니다. 친구가 넘어져서 아파할 때 "정말 많이 아프겠다. 괜찮니?"라고 말해주는 것이 바로 공감입니다. 반대로 "그걸 왜 못 보고 넘어지니?"라고 면박을 준다면 친구는 마음을 닫아버릴 것입니다.

공감하며 대화하려면 상대방의 눈을 바라보고 고개를 끄덕이며 반응을 보여주어야 합니다. 그리고 "그랬구나", "정말 힘들었겠다"처럼 맞장구를 쳐 주는 것도 좋습니다. 이렇게 내 마음을 알아주는 사람에게 우리는 진심을 털어놓게 됩니다. 좋은 대화는 화려한 말솜씨가 아니라, 상대방의 마음을 헤아리는 따뜻한 귀에서 시작됩니다.
\par\medskip\begin{itemize}[leftmargin=18pt, itemsep=2pt, topsep=2pt, label=\textbullet]
\item 공감(○): '정말 많이 아프겠다. 괜찮니?'
\item 면박(X): '그걸 왜 못 보고 넘어지니?'
\item 맞장구: '그랬구나', '정말 힘들었겠다'
\end{itemize}
}

\kfPassageNoteNT{%
중심 문장과 뒷받침 문장

* 중심 문장: 문단에서 가장 중요하고 중심이 되는 문장입니다. 주로 문단의 처음이나 끝에 위치하지만, 글의 구성에 따라 위치가 달라질 수 있습니다. * 뒷받침 문장: 중심 문장을 더 쉽게 이해할 수 있도록 예시나 까닭을 들어 자세히 설명하는 문장입니다.

문단을 읽을 때 전체 내용을 포괄하는 문장을 찾으면 중심 문장을, 이를 구체적으로 설명하는 문장을 찾으면 뒷받침 문장을 쉽게 파악할 수 있습니다.
\par\medskip\begin{itemize}[leftmargin=18pt, itemsep=2pt, topsep=2pt, label=\textbullet]
\item 중심 문장: '독서는 우리에게 많은 도움을 준다.'
\item 뒷받침 문장: '책을 읽으면 새로운 지식을 얻을 수 있다.'
\item 뒷받침 문장: '책을 읽으면 어휘력이 풍부해진다.'
\end{itemize}
}

\kfActivityBg \par\kfMaybeNeedspace{24\baselineskip}\noindent{\kfTipMark\,\,}{\small\color{kfInk} 다음 빈칸에 어울리는 어휘를 적어 보세요.}\par\nopagebreak[4]\smallskip\nopagebreak[4]
\begin{kfBlankList}
\kfBlankItem 친구의 고민을 들을 때는 딴짓을 하지 않고 \kfFillBlank[42mm]{}해야 한다.   (귀를 기울여 주의 깊게 듣는 것)
\kfBlankItem 내 말이 잘못 전달되어 친구와 작은 \kfFillBlank[42mm]{}가 생겼다.   (뜻을 잘못 알게 되는 것)
\end{kfBlankList}

\kfActivityBg \par\kfMaybeNeedspace{24\baselineskip}\noindent{\kfTipMark\,\,}{\small\color{kfInk} 다음 설명이 맞으면 ○, 틀리면 ×를 표시해 보세요.}\par\nopagebreak[4]\smallskip\nopagebreak[4]
\begin{kfOXList}
\kfOXItem 대화에서 가장 중요한 것은 내 말을 많이 하는 것이다. ( O / X ) \kfOX
\kfOXItem 상대방의 감정을 이해하고 자기도 그렇다고 느끼는 것을 '공감'이라고 한다. ( O / X ) \kfOX
\kfOXItem 상대방의 말에 고개를 끄덕이거나 맞장구를 치는 것은 좋은 태도이다. ( O / X ) \kfOX
\end{kfOXList}

\kfSecHeader{kfAreaConcept}{문제}{}

\begin{kfQBlock}{01}{객관식}
\kfQStem{\textbackslash{}<보기\textbackslash{}>에서 중심 문장은 몇 번 문장인가요?}
\begin{kfEvidence}① 건강을 지키기 위해 손 씻기는 매우 중요합니다.  \\\relax
 ② 손에는 눈에 보이지 않는 세균이 많기 때문입니다.  \\\relax
 ③ 외출 후에는 비누로 30초 이상 손을 씻어야 합니다.  \\\relax
④ 밥을 먹기 전에도 깨끗이 씻는 것이 좋습니다.\end{kfEvidence}
\begin{kfChoices}
\kfChoice ① 번 문장
\kfChoice ② 번 문장
\kfChoice ③ 번 문장
\kfChoice ④ 번 문장
\end{kfChoices}
\end{kfQBlock}
\begin{kfQBlock}{02}{객관식}
\kfQStem{\textbackslash{}<보기\textbackslash{}>에서 중심 문장은 몇 번 문장인가요?}
\begin{kfEvidence}① 횡단보도를 건널 때는 좌우를 잘 살펴야 합니다.  \\\relax
 ② 신호등이 초록불로 바뀌어도 바로 뛰지 말아야 합니다.  \\\relax
 ③ 길을 건너는 도중에는 장난을 치지 않습니다.  \\\relax
④ 이처럼 우리는 교통안전 수칙을 잘 지켜야 합니다.\end{kfEvidence}
\begin{kfChoices}
\kfChoice ① 번 문장
\kfChoice ② 번 문장
\kfChoice ③ 번 문장
\kfChoice ④ 번 문장
\end{kfChoices}
\end{kfQBlock}


\kfStartLit{「숲속 마을에 깃털이 알록달록한 앵무새 '…」}


\kfPassageNote{}{%
숲속 마을에 깃털이 알록달록한 앵무새 '초록이'가 살았습니다. 초록이는 다른 동물의 말을 흉내 내는 재주가 아주 뛰어났습니다. 하지만 말을 따라 하기만 할 뿐, 상대방의 기분은 전혀 생각하지 않는 버릇이 있었습니다.

어느 날 아기 다람쥐가 나무 아래서 엉엉 울고 있었습니다. 겨울 동안 먹으려고 모아 둔 도토리를 몽땅 잃어버렸기 때문입니다. 지나가던 곰 아저씨가 다람쥐를 다독이며 말했습니다.

"저런, 정말 속상하겠구나. 우리 같이 찾아보자."

그때 나무 위에 있던 초록이가 곰 아저씨의 말을 흉내 냈습니다.

"정말 속상하겠구나\textbackslash{}! 으하하, 속상하겠구나\textbackslash{}!"

초록이는 아주 신나고 장난스러운 목소리로 말을 따라 했습니다. 그 소리를 들은 아기 다람쥐는 자신의 슬픈 마음을 초록이가 놀린다고 생각했습니다. 다람쥐는 서러움이 북받쳐 더 크게 울음을 터뜨렸습니다. 곰 아저씨는 찌푸린 얼굴로 나무 위를 올려다보았습니다.

"초록아, 친구가 슬퍼할 때는 장난스럽게 말하면 안 된단다."

초록이는 고개를 갸웃거렸습니다. 자신은 그저 들리는 대로 따라 했을 뿐인데 왜 곰 아저씨가 화를 내는지 이해할 수 없었습니다. 그때 지혜로운 부엉이 할머니가 날아와 초록이 옆에 앉았습니다.

"초록아, 말은 입으로만 하는 게 아니란다. 마음을 담아서 해야 하지. 다람쥐가 지금 얼마나 슬플지 마음으로 먼저 느껴보렴."

초록이는 나무 아래서 눈물을 뚝뚝 흘리는 다람쥐를 가만히 내려다보았습니다. 그러자 가슴 한구석이 찌릿하며 미안한 마음이 들었습니다. 초록이는 나무 아래로 내려가 다람쥐에게 조그맣게 말했습니다.

"미안해, 다람쥐야. 많이 슬프지? 나도 같이 도토리를 찾아줄게."

초록이의 진심 어린 목소리에 아기 다람쥐는 울음을 그치고 고개를 끄덕였습니다. 그제야 초록이는 알게 되었습니다. 똑같은 말이라도 마음을 담지 않으면 상대방에게 상처가 될 수 있지만, 진심을 담으면 친구를 위로하는 따뜻한 선물이 된다는 것을 말입니다.
}


\kfActivityBg \par\kfMaybeNeedspace{18\baselineskip}\noindent{\kfTipMark\,\,}{\small\color{kfInk} 주어진 안내에 맞춰 자유롭게 글로 표현해 보세요.}\par\nopagebreak[4]\smallskip\nopagebreak[4]
\kfWritingPrompt{만약 여러분이 이야기 속의 초록이라면, 울고 있는 다람쥐에게 처음부터 어떻게 말해 주었을까요? 위로하는 말을 상상해서 적어보세요.}
\kfWriting{답안 작성}{8}

\kfSecHeader{kfAreaLiterature}{문제}{}

\begin{kfQBlock}{01}{객관식}
\kfQStem{아기 다람쥐가 나무 아래서 울고 있었던 원인은 무엇인가요?}
\begin{kfChoices}
\kfChoice 나무에서 떨어져서 다리를 다쳐서
\kfChoice 친구들과 술래잡기를 하다가 넘어져서
\kfChoice 앵무새 초록이가 깃털을 뽑아가서
\kfChoice 모아 둔 도토리를 몽땅 잃어버려서
\end{kfChoices}
\end{kfQBlock}
\begin{kfQBlock}{02}{객관식}
\kfQStem{부엉이 할머니가 초록이에게 가르쳐 준 '말하는 방법'은 무엇인가요?}
\begin{kfChoices}
\kfChoice 천둥처럼 목소리를 아주 크게 내는 것
\kfChoice 입으로만 하지 않고 마음을 담아서 하는 것
\kfChoice 어려운 단어를 골라서 많이 사용하는 것
\kfChoice 남의 말을 더 똑같이 흉내 내며 따라 하는 것
\end{kfChoices}
\end{kfQBlock}
\begin{kfQBlock}{03}{서술형}
\kfQStem{초록이가 곰 아저씨의 말을 따라 했을 때, 아기 다람쥐는 왜 더 크게 울었나요? 초록이가 자신의 슬픈 마음을 \_\_\_\_\_\_\_\_\_\_\_\_\_\_\_\_\_\_\_\_\_\_\_\_ 생각했기 때문에}
\kfAnswerNote{답안}{4}
\end{kfQBlock}
\begin{kfQBlock}{04}{객관식}
\kfQStem{초록이가 처음에 곰 아저씨가 화내는 이유를 이해하지 못한 까닭은 무엇인가요?}
\begin{kfChoices}
\kfChoice 귀가 잘 들리지 않아서
\kfChoice 상대방의 기분을 생각하지 않고 말만 따라 했기 때문에
\kfChoice 곰 아저씨의 말이 너무 빨라서
\kfChoice 다람쥐가 우는 소리가 너무 시끄러워서
\end{kfChoices}
\end{kfQBlock}
\begin{kfQBlock}{05}{서술형}
\kfQStem{이야기 끝부분에서 초록이가 "미안해, 많이 슬프지?"라고 말했을 때, 아기 다람쥐의 기분은 어떻게 변했을까요?}
\kfAnswerNote{답안}{4}
\end{kfQBlock}


\kfStartNonlit{공감하며 대화하는 방법}


\kfPassageNote{공감하며 대화하는 방법}{%
친구와 이야기를 나누는데 벽을 보고 말하는 것처럼 답답한 적이 있나요? 아니면 내 말에 친구가 고개를 끄덕여 주어서 마음이 뻥 뚫리는 것처럼 시원했던 적이 있나요? 대화는 혼자 하는 것이 아니라 두 사람이 함께 마음을 주고받는 것입니다. 그렇기 때문에 상대방의 마음을 얻는 '공감 대화법'이 필요합니다.

첫째, 귀로만 듣지 말고 눈과 몸으로 들어야 합니다. 이것을 '경청'이라고 합니다. 딴짓을 하며 건성으로 듣는 것은 상대방을 무시하는 행동입니다. 말하는 사람의 눈을 바라보고 몸을 그쪽으로 향하면, 상대방은 자신의 이야기가 존중받고 있다고 느낍니다. 또한 중간중간 고개를 끄덕여 주면 내 말을 잘 듣고 있다는 신호가 되어 상대방이 더 편안하게 이야기할 수 있습니다.

둘째, 적절한 '반응'을 보여 주어야 합니다. 이것을 맞장구라고도 합니다. 친구가 신나는 이야기를 할 때는 "정말?", "우와, 신기하다!"라며 같이 신나 하고, 슬픈 이야기를 할 때는 "저런, 속상했겠다"라며 낮은 목소리로 반응해 주는 것입니다. 아무 반응이 없으면 말하는 사람은 흥이 깨지거나 민망해질 수 있습니다.

셋째, 상대방의 감정을 읽어 주어야 합니다. 친구가 화를 낼 때 "너는 왜 맨날 화만 내니?"라고 따지기보다는 "지금 정말 화가 많이 났구나"라고 친구의 기분을 먼저 말로 표현해 주세요. 자신의 감정을 알아주는 사람에게 우리는 마음의 문을 열게 됩니다.

대화의 기술보다 더 중요한 것은 상대를 아끼는 진심입니다. 오늘부터 친구의 눈을 보고, 고개를 끄덕이며, 따뜻한 맞장구를 쳐 주는 공감 대장이 되어 보는 건 어떨까요?
}


\kfActivityBg \par\needspace{15\baselineskip}
\kfSecHeader{kfAreaNonlit}{활동}{문장 독해}
\par\kfMaybeNeedspace{32\baselineskip}\noindent{\kfTipMark\,\,}{\small\color{kfInk} 각 문장을 읽고 물음에 답하세요.}\par\nopagebreak[4]\smallskip\nopagebreak[4]
\begin{kfSentenceUnit}{1}{거친 파도가 모래사장으로 밀려온다}
\kfSentQ{1.}{(1) 이 문장에서 '무엇이'에 해당하는 것은 무엇인가요?}
\begin{kfChoices}
\kfChoice 거친
\kfChoice 파도가
\kfChoice 모래사장으로
\kfChoice 밀려온다
\end{kfChoices}
\kfSentQ{1.}{(2) 이 문장에서 '어찌하다'에 해당하는 것은 무엇인가요?}
\begin{kfChoices}
\kfChoice 거친
\kfChoice 파도가
\kfChoice 모래사장으로
\kfChoice 밀려온다
\end{kfChoices}
\end{kfSentenceUnit}

\begin{kfSentenceUnit}{2}{숲속의 다람쥐가 도토리를 먹는다}
\kfSentQ{2.}{(1) 이 문장에서 '무엇이'에 해당하는 것은 무엇인가요?}
\begin{kfChoices}
\kfChoice 숲속의
\kfChoice 다람쥐가
\kfChoice 도토리를
\kfChoice 먹는다
\end{kfChoices}
\kfSentQ{2.}{(2) 이 문장에서 '어찌하다'에 해당하는 것은 무엇인가요?}
\begin{kfChoices}
\kfChoice 숲속의
\kfChoice 다람쥐가
\kfChoice 맛있게
\kfChoice 먹는다
\end{kfChoices}
\end{kfSentenceUnit}

\kfActivityBg \kfSecHeader{kfAreaNonlit}{활동}{지시어 연결}
\par\kfMaybeNeedspace{18\baselineskip}\noindent{\kfTipMark\,\,}{\small\color{kfInk} 다음 문장에서 지시어가 가리키는 대상을 찾으세요.}\par\nopagebreak[4]\smallskip\nopagebreak[4]

\kfActivityBg \par\kfMaybeNeedspace{24\baselineskip}\noindent{\kfTipMark\,\,}{\small\color{kfInk} 빈칸에 알맞은 말을 채워 보세요.}\par\nopagebreak[4]\smallskip\nopagebreak[4]
\begin{kfBlankList}
\kfBlankItem 요약문 완성하기: 각 문단의 중심 내용을 잘 읽고, 빈칸에 들어갈 알맞은 말을 지문에서 찾아 쓰세요.

1문단: 대화는 두 사람이 함께 \kfFillBlank[42mm]{}을/를 주고받는 것이다. 2문단: 귀로만 듣지 말고 눈과 몸으로 듣는 \kfFillBlank[42mm]{}을/를 해야 한다. 3문단: 상대방의 말에 적절한 \kfFillBlank[42mm]{}을/를 보여 주어야 한다. 4문단: 상대방의 \kfFillBlank[42mm]{}을/를 읽어 주어야 마음의 문이 열린다. 5문단: 대화의 기술보다 더 중요한 것은 상대를 아끼는 \kfFillBlank[42mm]{}이다.
\end{kfBlankList}

\kfActivityBg \par\kfMaybeNeedspace{18\baselineskip}\noindent{\kfTipMark\,\,}{\small\color{kfInk} 주어진 안내에 맞춰 자유롭게 글로 표현해 보세요.}\par\nopagebreak[4]\smallskip\nopagebreak[4]
\kfWritingPrompt{지문에서 배운 대화 방법을 활용하여, 친구가 속상해할 때 따뜻한 맞장구와 공감의 말로 위로했던 경험(또는 앞으로 어떻게 위로할지)을 2\textasciitilde{}3문장으로 자세히 적어보세요.}
\kfWriting{답안 작성}{8}

\kfSecHeader{kfAreaNonlit}{문제}{}

\begin{kfQBlock}{01}{객관식}
\kfQStem{윗글의 '이것'이 가리키는 행동으로 가장 알맞은 것은?}
\begin{kfChoices}
\kfChoice 귀로만 듣는 것
\kfChoice 눈과 몸으로 듣는 것
\kfChoice 눈으로만 듣는 것
\kfChoice 다른 사람의 이야기를 무시하는 것
\end{kfChoices}
\end{kfQBlock}
\begin{kfQBlock}{02}{객관식}
\kfQStem{이 글에서 설명하는 '대화'의 의미로 가장 적절한 것은 무엇인가요?}
\begin{kfChoices}
\kfChoice 한 사람이 정해진 시간 동안 말을 끝까지 이어서 하는 것
\kfChoice 두 사람이 함께 마음을 주고받는 것
\kfChoice 듣는 사람이 말하는 사람의 의견에 무조건 동의해 주는 것
\kfChoice 상대의 말을 자르고 자신이 더 좋은 의견을 빨리 말하는 것
\end{kfChoices}
\end{kfQBlock}
\begin{kfQBlock}{03}{객관식}
\kfQStem{이 글에서 '경청'을 하는 방법으로 든 것이 \uline{\underline{아닌}} 것은 무엇인가요?}
\begin{kfChoices}
\kfChoice 말하는 사람의 눈을 바라본다.
\kfChoice 몸을 말하는 사람 쪽으로 향한다.
\kfChoice 중간중간 고개를 끄덕여 잘 듣고 있다는 신호를 보낸다.
\kfChoice 친구의 감정을 읽어 주는 말로 먼저 표현한다.
\end{kfChoices}
\end{kfQBlock}
\begin{kfQBlock}{04}{객관식}
\kfQStem{이 글에서 "벽을 보고 말하는 것처럼 답답하다"는 표현으로 비유한 상황은 무엇인가요?}
\begin{kfChoices}
\kfChoice 상대방이 내 말에 아무 반응을 보이지 않을 때
\kfChoice 친구가 "정말?", "우와, 신기하다!"라며 같이 신나 할 때
\kfChoice 친구가 고개를 끄덕여 주어 마음이 뻥 뚫리는 것처럼 시원할 때
\kfChoice 친구가 내 감정을 "지금 정말 화가 많이 났구나"라고 알아줄 때
\end{kfChoices}
\end{kfQBlock}
\begin{kfQBlock}{05}{객관식}
\kfQStem{이 글에서 알려 준 '셋째 방법(상대방의 감정 읽어 주기)'을 가장 잘 실천한 말은 무엇인가요?}
\begin{kfChoices}
\kfChoice (친구가 화를 낼 때) "너는 왜 맨날 화만 내니?"
\kfChoice (친구가 신나는 이야기를 할 때) "우와, 신기하다!"
\kfChoice (친구가 시험을 못 봐 우울해할 때) "시험을 망쳐서 정말 속상하겠구나. 기운 내."
\kfChoice (친구가 슬픈 이야기를 할 때) 고개를 끄덕이며 눈을 바라본다.
\end{kfChoices}
\end{kfQBlock}


\kfStartWeekly{실력확인 영역 학습}


\noindent\textit{\small 제한시간: 30분 / 총 12문항}\par\medskip
\par\noindent\textbf{[4\textasciitilde{}5] 다음 글을 읽고 물음에 답하시오.}\par\medskip
\kfPassageNote{지문 4}{%
앵무새의 말실수 숲속 마을에 깃털이 알록달록한 앵무새 '초록이'가 살았습니다. 초록이는 다른 동물의 말을 흉내 내는 재주가 아주 뛰어났습니다. 하지만 말을 따라 하기만 할 뿐, 상대방의 기분은 전혀 생각하지 않는 버릇이 있었습니다. 어느 날 아기 다람쥐가 나무 아래서 엉엉 울고 있었습니다. 겨울 동안 먹으려고 모아 둔 도토리를 몽땅 잃어버렸기 때문입니다. 지나가던 곰 아저씨가 다람쥐를 다독이며 말했습니다. "저런, 정말 속상하겠구나. 우리 같이 찾아보자." 그때 나무 위에 있던 초록이가 곰 아저씨의 말을 흉내 냈습니다. "정말 속상하겠구나! 으하하, 속상하겠구나!" 초록이는 아주 신나고 장난스러운 목소리로 말을 따라 했습니다. 그 소리를 들은 아기 다람쥐는 자신의 슬픈 마음을 초록이가 놀린다고 생각했습니다. 다람쥐는 서러움이 북받쳐 더 크게 울음을 터뜨렸습니다. 곰 아저씨는 찌푸린 얼굴로 나무 위를 올려다보았습니다. "초록아, 친구가 슬퍼할 때는 장난스럽게 말하면 안 된단다." 초록이는 고개를 갸웃거렸습니다. 자신은 그저 들리는 대로 따라 했을 뿐인데 왜 곰 아저씨가 화를 내는지 이해할 수 없었습니다. 그때 지혜로운 부엉이 할머니가 날아와 초록이 옆에 앉았습니다. "초록아, 말은 입으로만 하는 게 아니란다. 마음을 담아서 해야 하지. 다람쥐가 지금 얼마나 슬플지 마음으로 먼저 느껴보렴." 초록이는 나무 아래서 눈물을 뚝뚝 흘리는 다람쥐를 가만히 내려다보았습니다. 그러자 가슴 한구석이 찌릿하며 미안한 마음이 들었습니다. 초록이는 나무 아래로 내려가 다람쥐에게 조그맣게 말했습니다. "미안해, 다람쥐야. 많이 슬프지? 나도 같이 도토리를 찾아줄게." 초록이의 진심 어린 목소리에 아기 다람쥐는 울음을 그치고 고개를 끄덕였습니다. 그제야 초록이는 알게 되었습니다. 똑같은 말이라도 마음을 담지 않으면 소음에 불과하지만, 진심을 담으면 친구를 위로하는 따뜻한 선물이 된다는 것을 말입니다.
}
\par\noindent\textbf{[6\textasciitilde{}7] 다음 글을 읽고 물음에 답하시오.}\par\medskip
\kfPassageNote{지문 6}{%
공감하며 대화하는 방법 친구와 이야기를 나누는데 벽을 보고 말하는 것처럼 답답한 적이 있나요? 아니면 내 말에 친구가 고개를 끄덕여 주어서 마음이 뻥 뚫리는 것처럼 시원했던 적이 있나요? 대화는 혼자 하는 것이 아니라 두 사람이 함께 마음을 주고받는 것입니다. 그렇기 때문에 상대방의 마음을 얻는 '공감 대화법'이 필요합니다. 첫째, 귀로만 듣지 말고 눈과 몸으로 들어야 합니다. 이것을 '경청'이라고 합니다. 딴짓을 하며 건성으로 듣는 것은 상대방을 무시하는 행동입니다. 말하는 사람의 눈을 바라보고 몸을 그쪽으로 향하면, 상대방은 자신의 이야기가 존중받고 있다고 느낍니다. 또한 중간중간 고개를 끄덕여 주면 내 말을 잘 듣고 있다는 신호가 되어 상대방이 더 편안하게 이야기할 수 있습니다. 둘째, 적절한 '반응'을 보여 주어야 합니다. 이것을 맞장구라고도 합니다. 친구가 신나는 이야기를 할 때는 "정말?", "우와, 신기하다!"라며 같이 신나 하고, 슬픈 이야기를 할 때는 "저런, 속상했겠다"라며 낮은 목소리로 반응해 주는 것입니다. 아무 반응이 없으면 말하는 사람은 흥이 깨지거나 민망해질 수 있습니다. 셋째, 상대방의 감정을 읽어 주어야 합니다. 친구가 화를 낼 때 "너는 왜 맨날 화만 내니?"라고 따지기보다는 "지금 정말 화가 많이 났구나"라고 친구의 기분을 먼저 말로 표현해 주세요. 자신의 감정을 알아주는 사람에게 우리는 마음의 문을 열게 됩니다. 대화의 기술보다 더 중요한 것은 상대를 아끼는 진심입니다. 오늘부터 친구의 눈을 보고, 고개를 끄덕이며, 따뜻한 맞장구를 쳐 주는 공감 대장이 되어 보는 건 어떨까요?
}
\begin{kfQBlock}{01}{객관식}
\kfQStem{다음 설명에 알맞은 낱말은 무엇인가요?}
\begin{kfEvidence}"귀를 기울여 남의 말을 주의 깊게 듣는 것"\end{kfEvidence}
\begin{kfChoices}
\kfChoice 대화
\kfChoice 경청
\kfChoice 오해
\kfChoice 칭찬
\end{kfChoices}
\end{kfQBlock}

\begin{kfQBlock}{02}{객관식}
\kfQStem{다음 중 '진심'의 반대말로 가장 적절한 것은 무엇인가요?}
\begin{kfChoices}
\kfChoice 거짓
\kfChoice 감동
\kfChoice 사랑
\kfChoice 믿음
\end{kfChoices}
\end{kfQBlock}

\begin{kfQBlock}{03}{객관식}
\kfQStem{곰 아저씨가 초록이에게 화를 낸 이유는 무엇인가요?}
\begin{kfChoices}
\kfChoice 초록이가 노래를 너무 못 불러서
\kfChoice 초록이가 다람쥐의 도토리를 훔쳐 가서
\kfChoice 초록이가 슬퍼하는 친구에게 장난스럽게 말을 해서
\kfChoice 초록이가 나무 위에서 내려오지 않아서
\end{kfChoices}
\end{kfQBlock}

\begin{kfQBlock}{04}{객관식}
\kfQStem{이 글에서 친구가 내 말을 잘 듣고 있다는 신호가 되어 편안함을 주는 행동은 무엇인가요?}
\begin{kfChoices}
\kfChoice 눈을 감고 듣기
\kfChoice 창밖을 보며 듣기
\kfChoice 고개를 끄덕이기
\kfChoice 팔짱을 끼고 듣기
\end{kfChoices}
\end{kfQBlock}

\begin{kfQBlock}{05}{객관식}
\kfQStem{다음 문장의 짜임으로 알맞은 것은 무엇인가요?}
\begin{kfEvidence}"비행기가 하늘 높이 날아간다."\end{kfEvidence}
\begin{kfChoices}
\kfChoice 무엇이 무엇이다
\kfChoice 무엇이 어찌하다
\kfChoice 무엇이 어떠하다
\kfChoice 무엇이 누구이다
\end{kfChoices}
\end{kfQBlock}

\begin{kfQBlock}{06}{객관식}
\kfQStem{다음 문장의 짜임으로 알맞은 것은 무엇인가요?}
\begin{kfEvidence}"여름 날씨가 무척 덥다."\end{kfEvidence}
\begin{kfChoices}
\kfChoice 무엇이 무엇이다
\kfChoice 무엇이 어찌하다
\kfChoice 무엇이 어떠하다
\kfChoice 무엇이 언제이다
\end{kfChoices}
\end{kfQBlock}

\begin{kfQBlock}{07}{객관식}
\kfQStem{다음 중 주어와 서술어의 '시간 호응'이 바른 문장은 무엇인가요?}
\begin{kfChoices}
\kfChoice 나는 어제 학교에 갈 것이다.
\kfChoice 아까 점심을 맛있게 먹는다.
\kfChoice 내일 친구와 영화를 보았다.
\kfChoice 어제 밤에 무서운 꿈을 꾸었다.
\end{kfChoices}
\end{kfQBlock}

\begin{kfQBlock}{08}{서술형}
\kfQStem{초록이가 다람쥐를 내려다보며 미안한 마음이 생긴 까닭은 무엇인가요?}
\kfAnswerNote{답안}{4}
\end{kfQBlock}

\begin{kfQBlock}{09}{서술형}
\kfQStem{친구가 "나 어제 넘어져서 무릎이 너무 아파"라고 말했을 때, 가장 적절한 '반응'을 쓰시오.}
\kfAnswerNote{답안}{4}
\end{kfQBlock}

\begin{kfQBlock}{10}{서술형}
\kfQStem{다음 문장에서 어색한 부분을 바르게 고쳐 쓰시오.}
\begin{kfEvidence}"나는 내일 할머니 댁에 다녀왔다."\end{kfEvidence}
\kfAnswerNote{답안}{4}
\end{kfQBlock}

\begin{kfQBlock}{11}{서술형}
\kfQStem{이번 주에 배운 내용을 바탕으로, '친구와 대화를 잘하는 방법'을 3문장으로 서술하시오.}
\kfAnswerNote{답안}{4}
\end{kfQBlock}




\clearpage

\kfChapterCover{4}{소쉬르3 (챕터4)}{Chapter 4}{이번 챕터의 학습 내용을 시작합니다.}
\begin{kfChapterAreaList}
\kfChapterAreaItem{어휘}{kfAreaVocab}{생각의 깊이를 더하는 말들}
\kfChapterAreaItem{개념}{kfAreaConcept}{배경지식 넓히기: 내가 너라면 어떨까?}
\kfChapterAreaItem{문학}{kfAreaLiterature}{「오늘은 우리 둘이」}
\kfChapterAreaItem{비문학}{kfAreaNonlit}{다른 사람의 입장에서 생각하기}
\kfChapterAreaItem{실력확인}{kfAreaWeekly}{}
\end{kfChapterAreaList}

\kfStartVocab{생각의 깊이를 더하는 말들}


\begin{kfVocab}
\kfVocabItem{1}{처해 있는 사정이나 형편}
\kfVocabItem{2}{도와주거나 보살펴 주려는 마음}
\kfVocabItem{3}{사물이나 일을 바라보는 방향이나 생각}
\kfVocabItem{4}{높이어 귀중하게 대함}
\kfVocabItem{5}{일의 형편이나 그 일이 일어나게 된 까닭}
\kfVocabItem{6}{처지를 바꾸어 상대방의 입장에서 생각함}
\end{kfVocab}


\kfStartConcept{배경지식 넓히기: 내가 너라면 어떨까?}


\kfPassageNoteNT{%
역지사지의 마음

친구와 장난감을 가지고 싸운 적이 있나요? 나는 그냥 빌려달라고 한 것뿐인데 친구가 화를 내서 억울했을지도 모릅니다. 하지만 반대로 친구 입장에서는, 아끼는 장난감을 누가 함부로 만져서 속상했을 수도 있습니다. 누구나 자신만의 사정이 있기 마련입니다.

이렇게 자신의 생각만 고집하지 않고 상대방의 처지가 되어 생각해 보는 것을 사자성어로 '역지사지'라고 합니다. 이것은 서로의 입장을 바꾸어 생각한다는 뜻입니다. 내가 만약 선생님이라면 떠드는 학생을 볼 때 어떤 기분일지, 내가 만약 부모님이라면 밥을 안 먹는 자녀를 볼 때 얼마나 걱정이 될지 상상해 보는 것입니다.

사람마다 세상을 바라보는 관점은 모두 다릅니다. 내 눈에는 숫자 '6'으로 보이지만, 맞은편 친구의 관점에서는 '9'로 보일 수도 있습니다. 역지사지의 마음을 가지면 나와 다른 친구의 생각을 이해하고 존중하게 됩니다. 그리고 자연스럽게 상대를 위하는 배려의 마음도 자라납니다. 오늘 하루, 내 신발을 잠시 벗어두고 친구의 신발을 신은 것처럼 친구의 마음이 되어 생각해 보는 건 어떨까요?
\par\medskip\begin{itemize}[leftmargin=18pt, itemsep=2pt, topsep=2pt, label=\textbullet]
\item 장난감 다툼 — 친구의 입장(아끼는 장난감을 함부로 만져 속상)을 헤아리기
\item 내가 선생님이라면 — 떠드는 학생을 볼 때 어떤 기분일까?
\item 관점의 차이 — 내 눈의 '6'이 맞은편엔 '9'
\end{itemize}
}

\kfActivityBg \par\kfMaybeNeedspace{24\baselineskip}\noindent{\kfTipMark\,\,}{\small\color{kfInk} 다음 빈칸에 어울리는 어휘를 적어 보세요.}\par\nopagebreak[4]\smallskip\nopagebreak[4]
\begin{kfBlankList}
\kfBlankItem 친구와 다투었을 때는 서로의 \kfFillBlank[42mm]{}을 바꿔서 생각해 봐야 한다.   (처해 있는 형편)
\kfBlankItem 같은 그림이라도 보는 \kfFillBlank[42mm]{}에 따라 다르게 보일 수 있다.   (바라보는 방향이나 생각)
\end{kfBlankList}

\kfActivityBg \par\kfMaybeNeedspace{24\baselineskip}\noindent{\kfTipMark\,\,}{\small\color{kfInk} 다음 설명이 맞으면 ○, 틀리면 ×를 표시해 보세요.}\par\nopagebreak[4]\smallskip\nopagebreak[4]
\begin{kfOXList}
\kfOXItem '역지사지'는 처지를 바꾸어 상대방의 입장에서 생각하는 것이다. ( O / X ) \kfOX
\kfOXItem 사람마다 사물을 바라보는 관점은 모두 똑같다. ( O / X ) \kfOX
\kfOXItem 역지사지의 마음을 가지면 상대를 이해하고 존중하게 된다. ( O / X ) \kfOX
\end{kfOXList}


\kfStartLit{「오늘은 우리 둘이」}


\kfPassageNote{}{%
오늘은 우리 둘이 \\[4pt]
역할을 바꾸기로 했어요 \\[4pt]
나는 수염 난 아빠 \\[4pt]
아빠는 개구쟁이 아들 \\[14pt]
"아빠, 아이스크림 사 줘요." \\[4pt]
아빠가 조르자 내가 엄하게 말해요 \\[4pt]
"안 돼, 밥 다 먹고 먹어야지." \\[14pt]
"아빠, 텔레비전 더 볼래요." \\[4pt]
아빠가 떼를 쓰자 내가 짐짓 무섭게 말해요 \\[4pt]
"숙제는 다 했니? 일기 먼저 써야지." \\[14pt]
입을 삐죽거리는 아빠를 보니 \\[4pt]
웃음이 나면서도 알 것 같아요 \\[4pt]
매일매일 나를 따라다니던 \\[4pt]
아빠의 잔소리가 사랑이었다는 걸 \\[14pt]
아빠, 이제 알겠어요 \\[4pt]
아빠가 되는 건 참 힘든 일이에요
}


\kfActivityBg \par\kfMaybeNeedspace{18\baselineskip}\noindent{\kfTipMark\,\,}{\small\color{kfInk} 주어진 안내에 맞춰 자유롭게 글로 표현해 보세요.}\par\nopagebreak[4]\smallskip\nopagebreak[4]
\kfWritingPrompt{이번에는 여러분이 선생님이 되고, 학생이 되어보는 상황을 상상해 보세요. 선생님이 된 여러분은 학생에게 어떤 말을 해주고 싶나요?}
\kfWriting{답안 작성}{8}

\kfSecHeader{kfAreaLiterature}{문제}{}

\begin{kfQBlock}{01}{객관식}
\kfQStem{오늘 나와 아빠는 무엇을 하기로 했나요?}
\begin{kfChoices}
\kfChoice 서로 옷 바꿔 입기
\kfChoice 맛있는 요리 대결하기
\kfChoice 서로 역할 바꾸기
\kfChoice 숨바꼭질하기
\end{kfChoices}
\end{kfQBlock}
\begin{kfQBlock}{02}{객관식}
\kfQStem{화자가 역할을 바꾸어 본 뒤 깨달은 사실은 무엇인가요?}
\begin{kfChoices}
\kfChoice 아빠는 아이스크림을 정말 좋아한다.
\kfChoice 아빠의 잔소리는 나를 사랑하는 마음이었다.
\kfChoice 숙제를 하는 것은 정말 귀찮은 일이다.
\kfChoice 텔레비전은 밥을 먹고 봐야 재미있다.
\end{kfChoices}
\end{kfQBlock}
\begin{kfQBlock}{03}{서술형}
\kfQStem{아빠가 아이스크림을 사 달라고 조르자 '나'는 무엇이라고 말했나요?}
\kfAnswerNote{답안}{4}
\end{kfQBlock}
\begin{kfQBlock}{04}{객관식}
\kfQStem{시 속의 '나'가 아빠에게 "숙제는 다 했니?"라고 말했을 때, '나'의 마음은 어땠을까요?}
\begin{kfChoices}
\kfChoice 평소 아빠가 말씀하시는 말투를 따라 해 보고 싶은 호기심에서 한 말이었다.
\kfChoice 아빠가 훌륭한 사람이 되기를 바라는 걱정스러운 마음이었다.
\kfChoice 회사 일로 늦게 오신 아빠를 위로해 드리고 싶은 마음이었다.
\kfChoice 어른이 되어 봤더니 자녀에게는 잔소리할 필요가 없다고 알리고 싶었다.
\end{kfChoices}
\end{kfQBlock}
\begin{kfQBlock}{05}{객관식}
\kfQStem{마지막 연에서 "아빠가 되는 건 참 힘든 일이에요"라고 말한 까닭은 무엇일까요?}
\begin{kfChoices}
\kfChoice 회사 일이 많아서 자녀와 놀아 주기가 어려워서
\kfChoice 자녀에게 잔소리하면 자녀가 늘 화를 내서
\kfChoice 하고 싶은 것을 참으며 자녀를 바르게 가르쳐야 해서
\kfChoice 어른들은 무서운 표정만 지으며 살아야 하기 때문에
\end{kfChoices}
\end{kfQBlock}


\kfStartNonlit{다른 사람의 입장에서 생각하기}


\kfPassageNote{다른 사람의 입장에서 생각하기}{%
우리는 친구와 이야기하다가 서로 의견이 달라 다투기도 합니다. 나는 더운 날씨에 시원한 아이스크림을 먹고 싶은데, 친구는 배가 고프니 떡볶이를 먹자고 할 수도 있습니다. 이때 내 생각만 옳다고 고집하면 싸움이 일어납니다. 왜 우리는 서로 다른 생각을 할까요? 바로 각자 처해 있는 상황과 바라보는 관점이 다르기 때문입니다.

다툼을 줄이고 사이좋게 지내기 위해서는 '마법의 안경'이 필요합니다. 그것은 바로 '상대방의 입장이 되어 보는 안경'입니다. 이 안경을 쓰는 방법은 어렵지 않습니다.

첫째, "나라면 어땠을까?"라고 스스로에게 물어보는 것입니다. 약속 시간에 늦은 친구에게 화를 내기 전에, 내가 만약 급한 일이 생겨서 늦었다면 어떤 기분일지 상상해 보는 것입니다. 그러면 친구가 일부러 늦은 게 아니라 어쩔 수 없는 사정이 있었을 거라고 이해하게 됩니다.

둘째, 상대방의 이야기를 끝까지 들어주는 것입니다. 내 말만 앞세우면 상대방의 마음을 알 수 없습니다. 친구가 왜 떡볶이를 먹자고 했는지 이유를 들어보면, 아침을 굶어서 너무 배가 고팠다는 사실을 알게 될 수도 있습니다. 이유를 알게 되면 섭섭했던 마음이 사라지고 친구를 배려하고 싶은 마음이 생깁니다.

다른 사람의 입장에서 생각하는 것은 나를 낮추는 것이 아닙니다. 오히려 상대방을 존중하고 더 넓은 마음으로 세상을 바라보는 지혜로운 태도입니다. 오늘부터 내 마음속에 '상대방의 입장이 되어 보는 안경'을 하나씩 챙겨 두는 것은 어떨까요?
}


\kfActivityBg \par\needspace{15\baselineskip}
\kfSecHeader{kfAreaNonlit}{활동}{문장 독해}
\par\kfMaybeNeedspace{32\baselineskip}\noindent{\kfTipMark\,\,}{\small\color{kfInk} 각 문장을 읽고 물음에 답하세요.}\par\nopagebreak[4]\smallskip\nopagebreak[4]
\begin{kfSentenceUnit}{1}{마법의 안경을 다툼을 줄이기 위해 우리는 쓴다}
\kfSentQ{1.}{(1) 이 문장에서 '누가'에 해당하는 것은 무엇인가요?}
\begin{kfChoices}
\kfChoice 마법의 안경을
\kfChoice 다툼을 줄이기 위해
\kfChoice 우리는
\kfChoice 쓴다
\end{kfChoices}
\kfSentQ{1.}{(2) 이 문장에서 '무엇을'에 해당하는 것은 무엇인가요?}
\begin{kfChoices}
\kfChoice 다툼을 줄이기 위해
\kfChoice 마법의 안경을
\kfChoice 쓴다
\kfChoice 우리는
\end{kfChoices}
\kfSentQ{1.}{(3) 이 문장에서 '어찌하다'에 해당하는 것은 무엇인가요?}
\begin{kfChoices}
\kfChoice 마법의 안경을
\kfChoice 다툼을
\kfChoice 우리는
\kfChoice 쓴다
\end{kfChoices}
\end{kfSentenceUnit}

\begin{kfSentenceUnit}{2}{떡볶이를 원했다 친구는 배가 고파서}
\kfSentQ{2.}{(1) 이 문장에서 '누가'에 해당하는 것은 무엇인가요?}
\begin{kfChoices}
\kfChoice 떡볶이를
\kfChoice 원했다
\kfChoice 친구는
\kfChoice 배가 고파서
\end{kfChoices}
\kfSentQ{2.}{(2) 이 문장에서 '무엇을'에 해당하는 것은 무엇인가요?}
\begin{kfChoices}
\kfChoice 친구는
\kfChoice 떡볶이를
\kfChoice 원했다
\kfChoice 배가 고파서
\end{kfChoices}
\kfSentQ{2.}{(3) 이 문장에서 '어찌하다'에 해당하는 것은 무엇인가요?}
\begin{kfChoices}
\kfChoice 친구는
\kfChoice 배가
\kfChoice 떡볶이를
\kfChoice 원했다
\end{kfChoices}
\end{kfSentenceUnit}

\kfActivityBg \kfSecHeader{kfAreaNonlit}{활동}{지시어 연결}
\par\kfMaybeNeedspace{18\baselineskip}\noindent{\kfTipMark\,\,}{\small\color{kfInk} 다음 문장에서 지시어가 가리키는 대상을 찾으세요.}\par\nopagebreak[4]\smallskip\nopagebreak[4]

\kfActivityBg \par\kfMaybeNeedspace{24\baselineskip}\noindent{\kfTipMark\,\,}{\small\color{kfInk} 빈칸에 알맞은 말을 채워 보세요.}\par\nopagebreak[4]\smallskip\nopagebreak[4]
\begin{kfBlankList}
\kfBlankItem 요약문 완성하기: 각 문단의 중심 내용을 잘 읽고, 빈칸에 들어갈 알맞은 말을 지문에서 찾아 쓰세요.

1문단: 내 \kfFillBlank[42mm]{}만 옳다고 고집하면 싸움이 일어난다. 2문단: 다툼을 줄이려면 상대방의 \kfFillBlank[42mm]{}이 되어 보아야 한다. 3문단: 첫 번째 방법은 "\kfFillBlank[42mm]{} 어땠을까?"라고 스스로에게 물어보는 것이다. 4문단: 두 번째 방법은 상대방의 이야기를 \kfFillBlank[42mm]{} 들어주는 것이다. 5문단: 입장을 바꾸어 생각하는 것은 상대방을 \kfFillBlank[42mm]{}하는 지혜로운 태도이다.
\end{kfBlankList}

\kfActivityBg \par\kfMaybeNeedspace{18\baselineskip}\noindent{\kfTipMark\,\,}{\small\color{kfInk} 주어진 안내에 맞춰 자유롭게 글로 표현해 보세요.}\par\nopagebreak[4]\smallskip\nopagebreak[4]
\kfWritingPrompt{최근에 친구나 가족과 다투었거나 의견이 달랐던 경험을 떠올려 보세요. 그때 '마법의 안경'을 쓰고 상대방의 입장에서 생각했다면 상황이 어떻게 달라졌을지 2\textasciitilde{}3문장으로 자유롭게 글을 써보세요.}
\kfWriting{답안 작성}{8}

\kfSecHeader{kfAreaNonlit}{문제}{}

\begin{kfQBlock}{01}{객관식}
\kfQStem{윗글의 '그것'이 가리키는 대상으로 가장 알맞은 것은?}
\begin{kfChoices}
\kfChoice 다툼
\kfChoice 마법의 안경
\kfChoice 상대방
\kfChoice 사이좋게 지내는 것
\end{kfChoices}
\end{kfQBlock}
\begin{kfQBlock}{02}{객관식}
\kfQStem{이 글의 중심 내용으로 가장 적절한 것은 무엇인가요?}
\begin{kfChoices}
\kfChoice 다른 사람과 의견이 다를 때 자신의 생각을 분명하게 주장하는 방법
\kfChoice 친구의 이야기를 빠르게 듣고 더 좋은 답을 찾아 알려 주는 방법
\kfChoice 다른 사람의 입장에서 생각하는 것의 중요성과 방법
\kfChoice 같은 음식을 좋아하는 친구를 골라 함께 어울리는 방법
\end{kfChoices}
\end{kfQBlock}
\begin{kfQBlock}{03}{객관식}
\kfQStem{글쓴이는 상대방의 입장이 되어 보는 것을 무엇에 비유했나요?}
\begin{kfChoices}
\kfChoice 투명한 유리창
\kfChoice 마법의 안경
\kfChoice 시원한 아이스크림
\kfChoice 커다란 신발
\end{kfChoices}
\end{kfQBlock}
\begin{kfQBlock}{04}{서술형}
\kfQStem{상대방의 입장이 되어 보기 위해 스스로에게 던져야 할 질문은 무엇인가요? 본문에서 찾아 쓰세요.}
\kfAnswerNote{답안}{4}
\end{kfQBlock}
\begin{kfQBlock}{05}{서술형}
\kfQStem{글쓴이가 "내 말만 앞세우면 상대방의 마음을 알 수 없다"라고 말한 까닭은 무엇일까요?}
\kfAnswerNote{답안}{4}
\end{kfQBlock}
\begin{kfQBlock}{06}{서술형}
\kfQStem{다음 상황을 읽고 '마법의 안경'을 쓴 후 친구의 입장을 상상해 보세요.}
\begin{kfEvidence}친구가 청소 당번인데 청소를 하지 않고 집에 가버렸습니다. 처음에는 화가 났지만, '마법의 안경'을 쓰고 친구의 입장을 상상해 보세요. \\\relax
* 내 생각(처음): "청소도 안 하고 가다니 정말 무책임해\textbackslash{}!" \\\relax
* 친구의 입장 상상하기:\end{kfEvidence}
\kfAnswerNote{답안}{4}
\end{kfQBlock}


\kfStartWeekly{실력확인 영역 학습}


\noindent\textit{\small 제한시간: 30분 / 총 12문항}\par\medskip
\par\noindent\textbf{[4\textasciitilde{}5] 다음 글을 읽고 물음에 답하시오.}\par\medskip
\kfPassageNote{지문 4}{%
아빠와 아들 \\[14pt]
오늘은 우리 둘이 \\[4pt]
역할을 바꾸기로 했어요 \\[4pt]
나는 수염 난 아빠 \\[4pt]
아빠는 개구쟁이 아들 \\[14pt]
"아빠, 아이스크림 사 줘요." \\[4pt]
아빠가 조르자 내가 엄하게 말해요 \\[4pt]
"안 돼, 밥 다 먹고 먹어야지." \\[14pt]
"아빠, 텔레비전 더 볼래요." \\[4pt]
아빠가 떼를 쓰자 내가 짐짓 무섭게 말해요 \\[4pt]
"숙제는 다 했니? 일기 먼저 써야지." \\[14pt]
입을 삐죽거리는 아빠를 보니 \\[4pt]
웃음이 나면서도 알 것 같아요 \\[4pt]
매일매일 나를 따라다니던 \\[4pt]
아빠의 잔소리가 사랑이었다는 걸 \\[14pt]
아빠, 이제 알겠어요 \\[4pt]
아빠가 되는 건 참 힘든 일이에요
}
\par\noindent\textbf{[6\textasciitilde{}7] 다음 글을 읽고 물음에 답하시오.}\par\medskip
\kfPassageNote{지문 6}{%
다른 사람의 입장에서 생각하기 우리는 친구와 이야기하다가 서로 의견이 달라 다투기도 합니다. 나는 더운 날씨에 시원한 아이스크림을 먹고 싶은데, 친구는 배가 고프니 떡볶이를 먹자고 할 수도 있습니다. 이때 내 생각만 옳다고 고집하면 싸움이 일어납니다. 왜 우리는 서로 다른 생각을 할까요? 바로 각자 처해 있는 상황과 바라보는 관점이 다르기 때문입니다.

다툼을 줄이고 사이좋게 지내기 위해서는 '마법의 안경'이 필요합니다. 그것은 바로 '상대방의 입장이 되어 보는 안경'입니다. 이 안경을 쓰는 방법은 어렵지 않습니다.

첫째, "나라면 어땠을까?"라고 스스로에게 물어보는 것입니다. 약속 시간에 늦은 친구에게 화를 내기 전에, 내가 만약 급한 일이 생겨서 늦었다면 어떤 기분일지 상상해 보는 것입니다. 그러면 친구가 일부러 늦은 게 아니라 어쩔 수 없는 사정이 있었을 거라고 이해하게 됩니다.

둘째, 상대방의 이야기를 끝까지 들어주는 것입니다. 내 말만 앞세우면 상대방의 마음을 알 수 없습니다. 친구가 왜 떡볶이를 먹자고 했는지 이유를 들어보면, 아침을 굶어서 너무 배가 고팠다는 사실을 알게 될 수도 있습니다. 이유를 알게 되면 섭섭했던 마음이 사라지고 친구를 배려하고 싶은 마음이 생깁니다.

다른 사람의 입장에서 생각하는 것은 나를 낮추는 것이 아닙니다. 오히려 상대방을 존중하고 더 넓은 마음으로 세상을 바라보는 지혜로운 태도입니다. 오늘부터 내 마음속에 '상대방의 입장이 되어 보는 안경'을 하나씩 챙겨 두는 것은 어떨까요?
}
\begin{kfQBlock}{01}{객관식}
\kfQStem{다음 뜻풀이에 해당하는 사자성어는 무엇인가요?}
\begin{kfEvidence}"처지를 서로 바꾸어 상대방의 입장에서 생각함"\end{kfEvidence}
\begin{kfChoices}
\kfChoice 일거양득
\kfChoice 역지사지
\kfChoice 동문서답
\kfChoice 십시일반
\end{kfChoices}
\end{kfQBlock}

\begin{kfQBlock}{02}{객관식}
\kfQStem{다음 상황에 어울리는 마음가짐으로 가장 적절한 낱말은 무엇인가요?}
\begin{kfEvidence}"친구가 무거운 짐을 들고 갈 때, 내가 들어주겠다고 나서는 마음"\end{kfEvidence}
\begin{kfChoices}
\kfChoice 배려
\kfChoice 입장
\kfChoice 관점
\kfChoice 후회
\end{kfChoices}
\end{kfQBlock}

\begin{kfQBlock}{03}{객관식}
\kfQStem{이 시에서 '나'가 아빠(역할을 한 아들)에게 "안 돼"라고 말한 이유는 무엇인가요?}
\begin{kfChoices}
\kfChoice 아이스크림 가격이 너무 비쌌기 때문에
\kfChoice 아빠를 일부러 골탕 먹이고 싶어서
\kfChoice 밥을 다 먹고 나서 먹어야 하기 때문에
\kfChoice 가게에 아이스크림이 다 팔렸기 때문에
\end{kfChoices}
\end{kfQBlock}

\begin{kfQBlock}{04}{객관식}
\kfQStem{시의 끝부분에서 '나'는 아빠가 되는 것이 참 힘든 일이라고 느꼈습니다. 그 이유는 무엇인가요?}
\begin{kfChoices}
\kfChoice 가족을 위해 매일 밖에서 일해야 하기 때문에
\kfChoice 말을 안 듣는 아들을 무섭게 혼내야 하기 때문에
\kfChoice 하고 싶은 것을 참고 바르게 가르쳐야 하기 때문에
\kfChoice 아들의 투정을 받아주는 것이 너무 피곤하기 때문에
\end{kfChoices}
\end{kfQBlock}

\begin{kfQBlock}{05}{객관식}
\kfQStem{이 글에 따르면 친구와 다툼이 일어나는 까닭은 무엇인가요?}
\begin{kfChoices}
\kfChoice 서로 좋아하는 음식이 같아서
\kfChoice 서로의 상황과 관점이 달라서
\kfChoice 친구가 약속 시간을 잘 지켜서
\kfChoice '마법의 안경'을 쓰고 있어서
\end{kfChoices}
\end{kfQBlock}

\begin{kfQBlock}{06}{객관식}
\kfQStem{공식적인 상황에서 지켜야 할 예절로 알맞지 \underline{않은} 것은 무엇인가요?}
\begin{kfChoices}
\kfChoice 존댓말을 사용한다.
\kfChoice 친구에게 말하듯이 편하게 말한다.
\kfChoice 바른 자세로 또박또박 말한다.
\kfChoice 듣는 사람을 존중하는 마음을 가진다.
\end{kfChoices}
\end{kfQBlock}

\begin{kfQBlock}{07}{객관식}
\kfQStem{다음 중 높임 표현이 올바르게 쓰인 문장은 무엇인가요?}
\begin{kfChoices}
\kfChoice 할머니께서 밥을 먹습니다.
\kfChoice 아버지, 물 좀 줘요.
\kfChoice 선생님이 교실에 들어옵니다.
\kfChoice 어머니께서 편찮으셔서 누워 계십니다.
\end{kfChoices}
\end{kfQBlock}

\begin{kfQBlock}{08}{객관식}
\kfQStem{주어와 서술어의 호응이 자연스러운 문장은 무엇인가요?}
\begin{kfChoices}
\kfChoice 내 장래 희망은 축구를 좋아한다.
\kfChoice 내가 늦은 까닭은 늦잠을 잤기 때문이다.
\kfChoice 내 취미는 그림을 그리는 화가이다.
\kfChoice 우리가 해야 할 일은 쓰레기를 줍자.
\end{kfChoices}
\end{kfQBlock}

\begin{kfQBlock}{09}{서술형}
\kfQStem{사람마다 사물을 바라보는 방향이나 생각이 다를 수 있습니다. 이것을 무엇이라고 하나요?}
\kfAnswerNote{답안}{4}
\end{kfQBlock}

\begin{kfQBlock}{10}{서술형}
\kfQStem{글쓴이가 말한 '마법의 안경'은 무엇을 의미하나요?}
\kfAnswerNote{답안}{4}
\end{kfQBlock}

\begin{kfQBlock}{11}{서술형}
\kfQStem{다음 문장에서 잘못된 높임 표현을 찾아 바르게 고쳐 쓰시오.}
\begin{kfEvidence}"할아버지께서 방에서 자요."\end{kfEvidence}
\kfAnswerNote{답안}{4}
\end{kfQBlock}

\begin{kfQBlock}{12}{서술형}
\kfQStem{이번 주에 배운 내용을 바탕으로, 친구와 다투었을 때 화해하기 위해 어떻게 생각해야 하는지 3문장으로 서술하시오.}
\kfAnswerNote{답안}{4}
\end{kfQBlock}




\end{document}
